\subsection{Rigidity and a numerical criterion for the Mumford
classes}\label{ssec:mumfordrigid}

The first case of (F2) is a propagation problem along one compact Shimura
curve: the exceptional classes are algebraic at every CM point
(\Cref{thm:mumfordcm}), and they are algebraic on every member exactly when
$B(W\times_{C}W)$ holds (\Cref{cor:mumfordequiv}). We ask here whether the
propagation can be obtained without the Lefschetz standard conjecture, by one
of the mechanisms that have proved other cases of the Hodge conjecture:
degeneration, deformation into a larger family on which the class stays of
Hodge type, specialisation to special points, or semiregularity. The first
two cannot work; the third works exactly when the representing cycles have
bounded degree; and the fourth reduces to one dimension count at one CM
point. We do not prove that the classes are algebraic at a very general
point.

\begin{setupenv}[The classes $\omega_{a}$]\label{setup:omegaa}
Keep \Cref{setup:mumford} and let $W\to C$ be a family as in
\Cref{cor:mumfordlefschetz}, with fibres $X_{t}$. Over $\bar{\QQ}$ let
$\pi_{0},\pi_{12},\pi_{13},\pi_{23}\in H^{2}(X)\otimes H^{2}(X)\subset
H^{4}(X\times X)$ be the classes that correspond, under the isomorphism
$H^{2}\otimes H^{2}\cong\End H^{2}$ given by the form ${\textstyle\bigwedge^{2}}\psi$
on $H^{2}(X)={\textstyle\bigwedge^{2}}V^{\vee}$, to the projections onto $\QQ\theta$,
$U_{12}$, $U_{13}$ and $U_{23}$ of \Cref{prop:mumfordrm}(i), and for
$a=(a_{0},a_{1},a_{2},a_{3})$ put
\[
  \omega_{a}=a_{0}\pi_{0}+a_{1}\pi_{12}+a_{2}\pi_{13}+a_{3}\pi_{23},
  \qquad L_{0}(a)=a_{0}+a_{1}+a_{2}+a_{3}.
\]
The four tensors are invariant under $G$ and under the exchange $\tau$ of the
two factors of $X\times X$, because ${\textstyle\bigwedge^{2}}\psi$ is symmetric; so each
$\omega_{a}$ is a flat section over $C$ of Hodge type at every point. The
argument of \Cref{prop:mumfordrm}(ii) shows that the rational ones are those
with $a_{0}\in\QQ$ and $(a_{1},a_{2},a_{3})=(\sigma_{3}(\alpha),
\sigma_{2}(\alpha),\sigma_{1}(\alpha))$ for some $\alpha\in F$; we call such a
class \emph{exceptional} when $\alpha\notin\QQ$. The products of divisor
classes in $H^{2}\otimes H^{2}$ are the rational $\omega_{a}$ with
$\alpha\in\QQ$, and $\pi_{0}$ is a multiple of
$\mathrm{pr}_{1}^{*}\theta\cup\mathrm{pr}_{2}^{*}\theta$; so adding a rational
multiple of $\pi_{0}$ to an exceptional class changes $L_{0}$ and not the
question of algebraicity. For $c\in C$ write $Y=X_{c}\times X_{c}$, let
$HT^{2}(Y)=H^{0,2}(Y)\oplus H^{1}(T_{Y})\oplus H^{0}({\textstyle\bigwedge^{2}}T_{Y})$,
of dimension $28+64+28=120$, act on $H^{*}(Y,\CC)$ by contraction as in
\Cref{setup:criterion}, and let $\kappa\in H^{1}(T_{Y})$ be the Kodaira--Spencer class
at $c$ of the family $W\times_{C}W\to C$. Put
\[
  e(\omega)=\dim HT^{2}(Y)-\dim\Ann_{HT^{2}(Y)}(\omega)
  =120-\dim\Ann_{HT^{2}(Y)}(\omega).
\]
\end{setupenv}

\begin{lemma}[Linear relations among conjugates]\label{lem:conjugates}
Let $F$ be a cubic number field with embeddings $\sigma_{1},\sigma_{2},
\sigma_{3}$, let $\alpha\in F\setminus\QQ$, and let
$c_{0},c_{1},c_{2},c_{3},a_{0}\in\QQ$. The three conjugates
$\sigma_{i}(\alpha)$ are pairwise distinct and nonzero, and if
$c_{0}a_{0}+\sum_{i}c_{i}\sigma_{i}(\alpha)=0$ then $c_{1}=c_{2}=c_{3}$.
\end{lemma}

\begin{proof}
The minimal polynomial of $\alpha$ has degree three, since $[F:\QQ]$ is prime,
and is separable, so the conjugates are distinct; they are nonzero because
$\alpha\ne0$. Let $\Gamma$ be the Galois group of a normal closure; it
permutes the $\sigma_{i}$ transitively by $\gamma\circ\sigma_{i}=
\sigma_{\gamma(i)}$. Put $\beta=\alpha-\Tr_{F/\QQ}(\alpha)/3$, so that
$\beta\ne0$ has trace zero, and $r=-c_{0}a_{0}-(c_{1}+c_{2}+c_{3})
\Tr_{F/\QQ}(\alpha)/3\in\QQ$; the hypothesis reads $\sum_{i}c_{i}\sigma_{i}
(\beta)=r$. Applying every $\gamma\in\Gamma$ and averaging, each $\sigma_{j}$
occurs $|\Gamma|/3$ times in the place of each $\sigma_{i}$, so
$r=\frac13(\sum_{i}c_{i})\sum_{j}\sigma_{j}(\beta)=0$. The relations
$c\in\QQ^{3}$ with $\sum_{i}c_{i}\sigma_{i}(\beta)=0$ form a $\Gamma$-stable
subspace $N$ containing $(1,1,1)$. The plane $\{\sum c_{i}=0\}$ is an
irreducible representation of $\Gamma$ over $\QQ$: being transitive on three
letters, $\Gamma$ is the alternating or the full symmetric group, so it
contains the $3$-cycles, which have no rational eigenvector in that plane.
If $N$ were larger than the line of $(1,1,1)$, it would therefore contain the
plane and be all of $\QQ^{3}$, and $\sigma_{1}(\beta)=0$. Hence
$c\in\QQ(1,1,1)$.
\end{proof}

\begin{lemma}[One exceptional class gives the others]\label{lem:oneexceptional}
Let $t\in C$. If one rational exceptional class $\omega$ is algebraic on
$X_{t}\times X_{t}$, then every Hodge class in $H^{2}(X_{t})\otimes
H^{2}(X_{t})$ is algebraic; in particular so are the two exceptional classes
of \Cref{prop:mumfordproduct}.
\end{lemma}

\begin{proof}
At a CM point every Hodge class of $X_{t}\times X_{t}$ is algebraic
(\Cref{thm:mumfordcm}). Let $t$ be any other point, so that the Hodge group
of $X_{t}$ is $G$ (proof of \Cref{cor:mumfordequiv}), and write $X=X_{t}$.
For $x\in H^{2}(X)\otimes H^{2}(X)$ let $E_{x}$ be the endomorphism $y\mapsto x_{*}(\theta^{2}\cup y)$ of $H^{2}(X,\QQ)$, where
$x_{*}$ is the action of $x$ as a correspondence. By Poincar\'e duality and
the hard Lefschetz isomorphism $\theta^{2}\colon H^{2}\to H^{6}$, the map
$x\mapsto E_{x}$ is a $G$-equivariant isomorphism onto $\End H^{2}(X,\QQ)$,
so it carries the Hodge classes onto the Hodge endomorphisms, which are
$\QQ\times F$ by \Cref{prop:mumfordrm}(ii). It carries the products of
divisor classes, which are invariant under $\Sp(V,\psi)$, onto the
$\Sp(V,\psi)$-equivariant endomorphisms, which are $\QQ\times\QQ$; hence it
carries $\omega$ to an element $(r,\beta)$ with $\beta\notin\QQ$. The
algebraic Hodge classes form a subspace $A$ containing the products of divisor
classes, and its image under $x\mapsto E_{x}$ is closed under composition:
if $x_{1},x_{2}$ are algebraic, then the composite correspondence $x_{1}\circ\Delta_{*}(\theta^{2})\circ
x_{2}$ is algebraic, and so is its K\"unneth component in
$H^{2}\otimes H^{2}$, since the K\"unneth projectors of $X\times X$ are
polynomials in the graphs of $[2]\times1$ and $1\times[2]$, which act on
$H^{i}(X)\otimes H^{j}(X)$ by $2^{i}$ and $2^{j}$ (as in Step 1 of the proof
of \Cref{thm:lefschetzfamily}); that component $x$ has
$E_{x}=E_{x_{1}}E_{x_{2}}$. So the image of $A$ is a subalgebra of
$\QQ\times F$ containing $\QQ\times\QQ$ and $(r,\beta)$, hence
$\QQ\times\QQ[\beta]=\QQ\times F$, because $\beta$ generates the cubic
field $F$. Thus $A$ is the space of all Hodge classes in
$H^{2}(X)\otimes H^{2}(X)$.
\end{proof}

\begin{theorem}[Rigidity of the exceptional classes]\label{thm:mumfordrigid}
Let $c\in C$, $Y=X_{c}\times X_{c}$, and let $\omega=\omega_{a}$ be a
rational exceptional class.
\begin{enumerate}[label=\textup{(\roman*)},leftmargin=2.6em,itemsep=2pt]
\item $\Ann_{H^{1}(T_{Y})}(\omega)=\CC\kappa$ and
  $\Ann_{H^{0,2}(Y)}(\omega)=0$.
\item $\Ann_{H^{0}(\bigwedge^{2}T_{Y})}(\omega)=0$ if and only if
  $L_{0}(a)=a_{0}+\Tr_{F/\QQ}(\alpha)\ne0$. In particular
  $\Ann_{HT^{2}(Y)}(\omega)=\CC\kappa$ and $e(\omega)=119$ whenever
  $L_{0}(a)\ne0$, which can always be arranged by adding a rational multiple
  of $\pi_{0}$.
\item Let $\mathfrak{S}$ be the Siegel space of weight one Hodge structures on
  $H^{1}(Y,\QQ)$ polarised by $\psi\oplus\psi$, of dimension $36$. The germ at
  $Y$ of the locus in $\mathfrak{S}$ where $\omega$ stays of Hodge type is the
  germ of the curve $\{X_{t}\times X_{t}\}$ through $Y$.
\end{enumerate}
\end{theorem}

\begin{proof}
\emph{Notation.} Over $\CC$ write $V=V_{1}\otimes V_{2}\otimes V_{3}$ and
$\psi=\epsilon_{1}\otimes\epsilon_{2}\otimes\epsilon_{3}$, as in the proof of
\Cref{lem:mumfordmu}, and identify $V$ with $H^{1}(X_{c},\CC)$ and with its
dual through $\psi$. As at every point of $C$, the Hodge cocharacter lies in
the third factor: $V_{3}^{1,0}=\CC e$, $V_{3}^{0,1}=\CC f$, $V^{1,0}=V_{1}
\otimes V_{2}\otimes e$ and $V^{0,1}=V_{1}\otimes V_{2}\otimes f$. Let
$N\in\mathfrak{sl}(V_{3})$ be the nilpotent element with $Ne=f$, $Nf=0$. Then
$H^{1}(Y,\CC)=V_{(1)}\oplus V_{(2)}$, two copies of $V$, with the form
$\psi\oplus\psi$, and $\omega\in{\textstyle\bigwedge^{2}}V_{(1)}\otimes
{\textstyle\bigwedge^{2}}V_{(2)}$.

\emph{Step 1: operators.} For $J\subseteq\{1,2,3\}$ let $E_{J}\subset\End V$
be the tensor product of $\mathfrak{sl}(V_{i})$ for $i\in J$ and of $\CC\cdot1$
for $i\notin J$. Then $\End V=\bigoplus_{J}E_{J}$, and the $E_{J}$ are
irreducible and pairwise non-isomorphic under conjugation by $G$. The
$\psi$-adjoint of an element of $\mathfrak{sl}(V_{i})$ is its negative, so
$E_{J}$ consists of $\psi$-skew operators when $|J|$ is odd and of
$\psi$-symmetric ones when $|J|$ is even. The map $v\wedge w\mapsto
\psi(v,\cdot)w-\psi(w,\cdot)v$ identifies ${\textstyle\bigwedge^{2}}V$ with the
$\psi$-symmetric operators $E_{\emptyset}\oplus E_{12}\oplus E_{13}\oplus
E_{23}$, which are $\CC\theta$, $U_{12}$, $U_{13}$, $U_{23}$; a direct
expansion shows that ${\textstyle\bigwedge^{2}}\psi$ becomes the form
$\langle M,M'\rangle=\frac12\operatorname{tr}(MM')$, and that the derivation
action of $X\in\End V$ becomes $M\mapsto XM+MX^{\psi}$, that is
$D_{X}M=[X,M]$ if $X$ is $\psi$-skew and $D_{X}M=XM+MX$ if $X$ is
$\psi$-symmetric. In particular $D_{X}^{\dagger}=\mp D_{X}$ in the two cases,
$\dagger$ denoting the adjoint for $\langle\,,\rangle$. By the Cayley--Hamilton
theorem $XY=\frac12[X,Y]+\frac12\operatorname{tr}(XY)$ for
$X,Y\in\mathfrak{sl}(V_{i})$. Finally, by \Cref{setup:omegaa},
$\omega=\sum_{j}M_{j}\otimes P_{a}M^{j}$ for a basis $(M_{j})$ of
${\textstyle\bigwedge^{2}}V$ with dual basis $(M^{j})$, where $P_{a}$ acts on
$E_{\emptyset},E_{12},E_{13},E_{23}$ by $a_{0},a_{1},a_{2},a_{3}$. As $P_{a}$
commutes with $G$, it preserves the Hodge types.

The space ${\textstyle\bigwedge^{2}}V^{0,1}$ consists of the $\psi$-symmetric
operators that kill $V^{0,1}$ and take values in it, namely
$\mathfrak{sl}(V_{1})\otimes1\otimes N\oplus1\otimes\mathfrak{sl}(V_{2})
\otimes N\subset E_{13}\oplus E_{23}$, on which $P_{a}$ is $a_{2}$ and $a_{3}$.
The map $v\otimes w\mapsto\psi(v,\cdot)w$ identifies $V^{0,1}\otimes V^{0,1}$
with $\Hom(V^{1,0},V^{0,1})=\End V_{1}\otimes\End V_{2}\otimes\CC N=
\bigoplus_{J\ni3}E_{J}^{-}$, where $E_{J}^{-}=E_{J}\cap\Hom(V^{1,0},V^{0,1})$
has dimension $1,3,3,9$ for $J=\{3\},\{1,3\},\{2,3\},\{1,2,3\}$.

\emph{Step 2: the pieces.} The space $HT^{2}(Y)$ acts as follows:
$H^{0,2}={\textstyle\bigwedge^{2}}H^{0,1}$ by the wedge product,
$H^{1}(T_{Y})=\Hom(H^{1,0},H^{0,1})$ by derivations, and
$H^{0}({\textstyle\bigwedge^{2}}T_{Y})={\textstyle\bigwedge^{2}}(H^{1,0})^{\vee}$,
which $\psi\oplus\psi$ identifies with ${\textstyle\bigwedge^{2}}H^{0,1}$, by
contraction. Write $\eta\in{\textstyle\bigwedge^{2}}H^{0,1}$ as
$\eta_{11}+\eta_{12}+\eta_{22}$ with $\eta_{11}\in{\textstyle\bigwedge^{2}}
V^{0,1}_{(1)}$, $\eta_{12}\in V^{0,1}_{(1)}\otimes V^{0,1}_{(2)}$,
$\eta_{22}\in{\textstyle\bigwedge^{2}}V^{0,1}_{(2)}$, and $\xi\in H^{1}(T_{Y})$
as a matrix $(\xi_{st})$ with $\xi_{st}\in\Hom(V^{1,0}_{(t)},V^{0,1}_{(s)})$.
The products $\eta_{11}\wedge\omega$, $\eta_{12}\wedge\omega$,
$\eta_{22}\wedge\omega$ lie in different summands
${\textstyle\bigwedge^{p}}V_{(1)}\otimes{\textstyle\bigwedge^{q}}V_{(2)}$, and so
do the three contractions, and $\xi_{11}\omega+\xi_{22}\omega$, $\xi_{12}\omega$,
$\xi_{21}\omega$. So each annihilator is computed piece by piece, and $\tau$
exchanges the pieces $11$ and $22$, $12$ and $21$. For a rational exceptional
class, \Cref{lem:conjugates} gives: $a_{1},a_{2},a_{3}$ are pairwise distinct
and nonzero, and no linear form in $a_{0},\dots,a_{3}$ with rational
coefficients that are not all equal on $a_{1},a_{2},a_{3}$ vanishes. We use
this for $a_{i}\pm a_{j}$ ($i\ne j$), $a_{2}+3a_{3}$, $3a_{2}+a_{3}$ and the
three forms in Step 4.

\emph{Step 3: $H^{0,2}$.} Let $(x_{j})$ be a basis of
${\textstyle\bigwedge^{2}}V^{1,0}$ and $(x^{j})$ the dual basis of
${\textstyle\bigwedge^{2}}V^{0,1}$. The component of $\omega$ in
${\textstyle\bigwedge^{2}}V^{1,0}_{(1)}\otimes{\textstyle\bigwedge^{2}}
V^{0,1}_{(2)}$ is $\sum_{j}x_{j}\otimes Kx^{j}$, where $K=P_{a}$ on
${\textstyle\bigwedge^{2}}V^{0,1}$ is invertible because $a_{2}a_{3}\ne0$. For
$\eta_{11}\ne0$ the part of $\eta_{11}\wedge\omega$ in
$H^{2,2}(X_{c})\otimes H^{0,2}(X_{c})$ is $\sum_{j}(\eta_{11}\wedge x_{j})
\otimes Kx^{j}$, which is nonzero because ${\textstyle\bigwedge^{2}}V^{0,1}
\otimes{\textstyle\bigwedge^{2}}V^{1,0}\to{\textstyle\bigwedge^{4}}V$ is
injective; by $\tau$ the same holds for $\eta_{22}$. For
$\eta_{12}=\sum_{r}v_{r}\otimes w_{r}$ the part of $\eta_{12}\wedge\omega$ in
$H^{2,1}\otimes H^{0,3}$ is $\sum_{r,j}(v_{r}\wedge x_{j})\otimes
(w_{r}\wedge Kx^{j})$. If it vanishes, then, as $V^{0,1}\otimes
{\textstyle\bigwedge^{2}}V^{1,0}\to{\textstyle\bigwedge^{3}}V$ is injective and
the $Kx^{j}$ span ${\textstyle\bigwedge^{2}}V^{0,1}$,
$\sum_{r}\phi(v_{r})\,w_{r}\wedge y=0$ for every linear form $\phi$ on
$V^{0,1}$ and every $y\in{\textstyle\bigwedge^{2}}V^{0,1}$; as
$\dim V^{0,1}=4$ this forces $\sum_{r}\phi(v_{r})w_{r}=0$ for every $\phi$,
so $\eta_{12}=0$. Hence $\Ann_{H^{0,2}(Y)}(\omega)=0$.

\emph{Step 4: $H^{0}({\textstyle\bigwedge^{2}}T_{Y})$.} Up to signs and one
nonzero factor fixed by the normalisation of the contraction,
$\iota_{\eta_{11}}\omega=\sum_{j}\langle\eta_{11},M_{j}\rangle P_{a}M^{j}=
P_{a}\eta_{11}$, which is nonzero for $\eta_{11}\ne0$ because $a_{2}a_{3}\ne0$;
likewise for $\eta_{22}$. For $\eta_{12}=v\otimes w$ the contraction is
$\sum_{j}M_{j}v\otimes(P_{a}M^{j})w$, and through $v\otimes w\mapsto
\psi(v,\cdot)w$ it becomes the operator $Q_{a}(R)=\sum_{j}(P_{a}M^{j})\,R\,
M_{j}$ with $R=\psi(v,\cdot)w$, because the $M_{j}$ are $\psi$-symmetric.
For a basis $(X_{k})$ of $\mathfrak{sl}(V_{i})$ and the dual basis $(X^{k})$
for the trace form, $\sum_{k}X^{k}RX_{k}$ is $\frac32R$ for $R=1$ and
$-\frac12R$ for $R\in\mathfrak{sl}(V_{i})$ (take $X_{k}=h,e,f$ and $X^{k}=h/2,
f,e$). The dual basis of $E_{ij}$ for $\langle\,,\rangle$ is the tensor
product of these dual bases, and that of $E_{\emptyset}$ is $\frac14$. Hence
$Q_{a}$ acts on $E_{J}$ by the scalar
\[
  q_{J}(a)=\tfrac14a_{0}+\sum_{\{i,j\}}a_{\{i,j\}}\,\lambda_{i}\lambda_{j},
  \qquad\lambda_{i}=-\tfrac12\ (i\in J),\quad\lambda_{i}=\tfrac32\
  (i\notin J),
\]
where $a_{\{1,2\}}=a_{1}$, $a_{\{1,3\}}=a_{2}$, $a_{\{2,3\}}=a_{3}$; on the four
pieces $E_{J}^{-}$ this is
\begin{gather*}
  4q_{3}=a_{0}+9a_{1}-3a_{2}-3a_{3},\qquad
  4q_{13}=a_{0}-3a_{1}+a_{2}-3a_{3},\\
  4q_{23}=a_{0}-3a_{1}-3a_{2}+a_{3},\qquad
  4q_{123}=L_{0}(a).
\end{gather*}
The first three do not vanish. So the annihilator in
$H^{0}({\textstyle\bigwedge^{2}}T_{Y})$ is zero when $L_{0}(a)\ne0$, and is the
piece $E_{123}^{-}$ of $\eta_{12}$, of dimension nine, when $L_{0}(a)=0$.

\emph{Step 5: $H^{1}(T_{Y})$, the diagonal pieces.} Put
$\Phi(x\otimes y)(z)=\langle z,x\rangle y$, so that $\Phi(\omega)=P_{a}$,
$\Phi((D\otimes1)\omega)=P_{a}D^{\dagger}$ and $\Phi((1\otimes D)\omega)=
DP_{a}$. Thus $\xi_{11}\omega+\xi_{22}\omega=0$ if and only if
$P_{a}D_{\xi_{11}}^{\dagger}+D_{\xi_{22}}P_{a}=0$. This condition is
$G$-equivariant in $(\xi_{11},\xi_{22})\in\End V\oplus\End V$, so its solutions
are the pairs $(b_{11}X,b_{22}X)$ with $X\in E_{J}$ and $(b_{11},b_{22})$ in a
subspace $M_{J}\subseteq\CC^{2}$; we need $J\ni3$. For $X=\bigotimes X_{i}\in
E_{J}$ and $M=\bigotimes M_{i}\in E_{I}$, the product rule above gives
$XM\pm MX$ as $2^{1-|I\cap J|}$ times the sum, over $T\subseteq I\cap J$ with
$|T|\equiv|J|$ modulo $2$, of the tensor with factor $[X_{i},M_{i}]$ for
$i\in T$, $\operatorname{tr}(X_{i}M_{i})$ for $i\in(I\cap J)\setminus T$, and
the remaining factor of $X$ or $M$ elsewhere; this term lies in
$E_{(I\triangle J)\cup T}$ and is nonzero for suitable $X$ and $M$. On the
block from $E_{I}$ to $E_{I'}$ the condition reads $b_{22}a_{I}=b_{11}a_{I'}$
for $|J|$ odd and $b_{22}a_{I}=-b_{11}a_{I'}$ for $|J|$ even, with
$a_{\emptyset}=a_{0}$, $a_{12}=a_{1}$, $a_{13}=a_{2}$, $a_{23}=a_{3}$.
\begin{itemize}[leftmargin=1.6em,itemsep=1pt]
\item $J=\{3\}$: the blocks from $E_{13}$ to $E_{13}$ and from $E_{23}$ to
  $E_{23}$ give $(b_{22}-b_{11})a_{2}=0$, so $M_{J}=\CC(1,1)$.
\item $J=\{1,2,3\}$: there is a block between any two distinct
  $E_{12},E_{13},E_{23}$, so $b_{22}a_{1}=b_{11}a_{2}$ and
  $b_{22}a_{2}=b_{11}a_{1}$, whence $(b_{11}+b_{22})(a_{1}-a_{2})=0$ and
  $b_{11}=-b_{22}$; then $b_{22}(a_{1}+a_{2})=b_{22}(a_{1}+a_{3})=0$, and
  $b_{22}\ne0$ would give $a_{2}=-a_{1}=a_{3}$. So $M_{J}=0$.
\item $J=\{1,3\}$: the blocks between $E_{12}$ and $E_{23}$ give
  $b_{22}a_{1}+b_{11}a_{3}=b_{22}a_{3}+b_{11}a_{1}=0$, of determinant
  $a_{1}^{2}-a_{3}^{2}\ne0$, so $M_{J}=0$; and $J=\{2,3\}$ likewise, with
  $a_{1}^{2}-a_{2}^{2}\ne0$.
\end{itemize}
So on the diagonal pieces the annihilator is the line of $(N,N)$, with $N$
acting on both copies of $V$.

\emph{Step 6: $H^{1}(T_{Y})$, the mixed pieces.} The map $\xi_{12}$ kills
$V^{0,1}_{(2)}$, so the part of $\xi_{12}\omega$ in $H^{0,3}(X_{c})\otimes
H^{1,0}(X_{c})$ comes from the component $\sum_{j}Kx^{j}\otimes x_{j}$ of
$\omega$ in ${\textstyle\bigwedge^{2}}V^{0,1}_{(1)}\otimes{\textstyle\bigwedge^{2}}
V^{1,0}_{(2)}$ only, and is $\sum_{j}Kx^{j}\wedge\xi_{12}(x_{j}')\otimes
x_{j}''$, where $x\mapsto x'\otimes x''$ is the coproduct
${\textstyle\bigwedge^{2}}\to V\otimes V$. Identify $V^{1,0}$ and $V^{0,1}$
with $W=V_{1}\otimes V_{2}$ through $e$ and $f$, with the symmetric form
$q=\epsilon_{1}\otimes\epsilon_{2}$ and a volume form; then
${\textstyle\bigwedge^{2}}W$ is the space of $q$-skew operators
$\mathfrak{sl}(V_{1})\otimes1\oplus1\otimes\mathfrak{sl}(V_{2})$, the Hodge
star is $+1$ on one summand and $-1$ on the other, and $K$ is $a_{2}$ on the
first and $a_{3}$ on the second. Pair the part above with $b\in W$ in the
first factor, through the volume form, and with $\beta\in W$ in the second,
through $q$. Writing $\operatorname{vol}(y\wedge y')=q(\ast y,y')$ on
${\textstyle\bigwedge^{2}}W$ turns the pairing into the trace of
$x\mapsto K\ast\bigl(\xi_{12}(\iota_{\beta}x)\wedge b\bigr)$ on
${\textstyle\bigwedge^{2}}W$, and summing over orthogonal bases of the two
summands gives, up to a nonzero factor and a sign,
$q\bigl(c_{a}(\xi_{12})\beta,b\bigr)$ with
\[
  c_{a}(\xi)=a_{2}\sum_{k}(X^{k}\otimes1)\,\xi\,(X_{k}\otimes1)
  -a_{3}\sum_{k}(1\otimes X^{k})\,\xi\,(1\otimes X_{k})
\]
for $\xi\in\End W=\End V_{1}\otimes\End V_{2}$.
By the scalars of Step 4, $c_{a}$ acts on $1\otimes1$, $\mathfrak{sl}(V_{1})
\otimes1$, $1\otimes\mathfrak{sl}(V_{2})$, $\mathfrak{sl}(V_{1})\otimes
\mathfrak{sl}(V_{2})$, that is on $E_{3}^{-}$, $E_{13}^{-}$, $E_{23}^{-}$,
$E_{123}^{-}$, by
\[
  \tfrac32(a_{2}-a_{3}),\qquad
  -\tfrac12(a_{2}+3a_{3}),\qquad
  \tfrac12(3a_{2}+a_{3}),\qquad
  -\tfrac12(a_{2}-a_{3}),
\]
none of which vanishes. So $\xi_{12}\omega\ne0$ for $\xi_{12}\ne0$, and by
$\tau$ the same holds for $\xi_{21}$.

\emph{Step 7: $\kappa$.} The family of Hodge structures along $C$ is the
orbit of the Hodge structure at $c$ under the real points of the third factor
of $G$, so the Kodaira--Spencer class $\kappa$ is the component of type
$(-1,1)$ of an element of $\mathfrak{sl}(V_{3})$ acting on both copies of $V$,
a nonzero multiple of $(N,N)$; it is nonzero because the Hodge structures
along $C$ vary. By Steps 5 and 6, $\Ann_{H^{1}(T_{Y})}(\omega)=\CC\kappa$.
With Steps 3 and 4 this proves (i) and (ii).

For (iii), a point of $\mathfrak{S}$ is a Hodge filtration $F^{1}$ on
$H^{1}(Y,\CC)$, and $\omega$ is of Hodge type at it when $\omega\in F^{2}
H^{4}$, a holomorphic condition. For weight one every tangent direction of
$\mathfrak{S}$ at $Y$ is an element $\xi$ of
$\mathfrak{sp}(\psi\oplus\psi)^{-1,1}\subset H^{1}(T_{Y})$, and it moves
$F^{2}H^{4}$ by the derivation $\xi$ of ${\textstyle\bigwedge^{*}}H^{1}$, which maps
$H^{2,2}$ into $H^{1,3}$; so the Zariski tangent space of the Hodge locus is
$\{\xi:\xi\lrcorner\,\omega=0\}$, which is $\CC\kappa$ by (i). This is the
usual description of the tangent space of the Hodge locus of a class of type
$(p,p)$ as the kernel of the contraction into $H^{p-1,p+1}$
\cite[Chapter~5]{Voi03}, in its simplest case: a variation of weight one is
subject to no horizontality condition, so every direction of $\mathfrak{S}$
enters. The Hodge locus is an analytic germ whose
Zariski tangent space is a line, so it lies in a smooth curve germ; it
contains the smooth curve germ $\{X_{t}\times X_{t}\}$, whose tangent line is
$\CC\kappa$; so the two germs are equal.
\end{proof}

\begin{remark}[Comparison with the divisor classes]\label{rem:rigidcompare}
For $\alpha\in\QQ$ the class $\omega_{a}$ is a combination of products of
divisor classes and is invariant under the whole of $\Sp(V,\psi)$ acting
diagonally on $V\oplus V$, so its Hodge locus in $\mathfrak{S}$ contains the
diagonal Siegel space $\{X'\times X'\}$, of dimension $10$; in the proof above
this shows in the scalars $a_{2}-a_{3}$ of Step~6, which vanish when
$a_{1}=a_{2}=a_{3}$. The products of divisor classes therefore move with
$X'\times X'$ for every $X'$, whereas the exceptional classes move with
$X_{t}\times X_{t}$ for $t\in C$ only and cease to be of Hodge type under
every deformation of $Y$ off the curve, however small.
\end{remark}

\begin{figure}[tbp]
\centering
  \includegraphics[width=\textwidth]{fig_rigidity}
\caption{\Cref{thm:mumfordrigid}. Top: the tangent space
$\mathfrak{sp}^{-1,1}$ of the Siegel space $\mathfrak{S}$ at
$Y=X_{c}\times X_{c}$, of dimension $36$, split into nine blocks by the
weights of a maximal torus of $\SL(V_{1})\times\SL(V_{2})$, one cube for each
dimension. The tangent space of the Hodge locus of an exceptional class is
the line $\CC\kappa$ of the Mumford curve, in the block of weight $(0,0)$;
for $\alpha\in\QQ$ it contains the tangent space of the diagonal Siegel space,
of dimension $10$ (\Cref{rem:rigidcompare}). Bottom: the scalars of Steps 4
to 6 of the proof on the pieces $E_{J}^{-}$ of $HT^{2}(Y)$; the pieces of
$H^{0,2}$ and the unmixed pieces of $H^{0}(\bigwedge^{2}T_{Y})$ are controlled
by the block $K$ of $P_{a}$ (Steps 3 and 4).}
\label{fig:rigidity}
\end{figure}

\begin{corollary}[No degeneration and no larger family]
\label{cor:nodegeneration}
Let $\omega$ be a rational exceptional class on $X_{c}\times X_{c}$. Every
connected complex analytic family of polarised abelian eightfolds through
$X_{c}\times X_{c}$ on which $\omega$ stays of Hodge type lies, near each
point of the curve, inside the curve $\{X_{t}\times X_{t}:t\in C\}$, which is
compact. In particular no such family degenerates, and on no member of it
other than the CM points of the curve is $\omega$ a combination of products
of divisor classes or of Weil classes of CM fields acting on the member or on
its quotients.
\end{corollary}

\begin{proof}
By \Cref{thm:mumfordrigid}(iii) the Hodge locus of $\omega$ coincides with the
curve in a neighbourhood of every point of the curve, so a connected family
through $Y$ on which $\omega$ stays of Hodge type maps into the curve near
each of its points over the curve; by connectedness and the identity theorem
its image lies in the curve. The curve is compact because $D$ is a division
algebra, so that the arithmetic group of the Shimura curve is cocompact
\cite{Mum69}. Indeed, by \Cref{setup:mumford} the algebra $D$ is definite at
$\sigma_{1}$ and $\sigma_{2}$, so $D\otimes_{F,\sigma_{1}}\RR$ is Hamilton's
quaternions and $D$ is not $M_{2}(F)$: it is a division algebra, split at the
one real place $\sigma_{3}$. The norm one group of an order in $D$, embedded
in $\SL_{2}(\RR)$ through $\sigma_{3}$, is a Fuchsian group of finite
covolume with no parabolic element, since a parabolic $\gamma\ne\pm1$ would
make $\gamma\mp1$ a nonzero nilpotent element of $D$; and a Fuchsian group of
finite covolume without parabolic elements has compact quotient.

For the last assertion: a member of the curve other than a CM point has
Hodge group $G$ (proof of \Cref{cor:mumfordequiv}), and on it $\omega$ lies
outside the span of the products of divisor classes and of the Weil classes
of quotients, by \Cref{prop:mumfordproduct}(ii) and (iii). At a CM point the
class is already algebraic (\Cref{thm:mumfordcm}), and the CM points, being
countably many, do not carry algebraicity to the other members
(\Cref{lem:spreading} needs uncountably many points).
\end{proof}

The same annihilator controls deformations that are not algebraic at all,
and so decides whether the twistor lines along which Markman moves sheaves
in his proof for Weil classes can be used here.

\begin{proposition}[The Hodge locus among all complex tori, and twistor
lines]\label{prop:notwistor}
Let $\omega$ be a rational exceptional class, $Y=X_{c}\times X_{c}$ with
$c\in C$, and $\Lambda=H_{1}(Y,\ZZ)$. Let $\mathcal{T}$ be the space of
complex structures on the real torus $\Lambda_{\RR}/\Lambda$, a complex
manifold of dimension $64$ with tangent space $H^{1}(T_{Y})$ at $Y$, and
$\mathcal{T}_{\omega}\subset\mathcal{T}$ the closed analytic subset on which
$\omega$ is of type $(2,2)$.
\begin{enumerate}[label=\textup{(\roman*)},leftmargin=2.6em,itemsep=2pt]
\item The connected component of $\mathcal{T}_{\omega}$ that contains $Y$ is
  the universal cover $\mathfrak{H}$ of the Mumford curve, the set of
  squares of members of the family, biholomorphic to the upper half plane.
  Every connected complex analytic family of complex tori, algebraic or not,
  through a member of the Mumford family on which $\omega$ stays of Hodge
  type lies in the family.
\item No twistor line through a member of the Mumford family, and no compact
  connected complex curve in $\mathcal{T}$ through it, keeps $\omega$ of
  Hodge type.
\item Nevertheless $\mathcal{T}_{\omega}$ contains twistor lines. The algebra
  $D\otimes_{F,\sigma_{1}}\RR$ is Hamilton's quaternions and acts on
  $V_{\RR}$ through the first factor of $G(\RR)$; for a unit quaternion $j$
  with $j^{2}=-1$, the complex structure $j\oplus j$ on $V_{\RR}\oplus
  V_{\RR}$ makes $\omega$ of type $(2,2)$, and these structures form a twistor
  line. On each factor $\psi(x,jy)$ has signature $(4,4)$, and a member whose
  Mumford--Tate group is $G$ is not an abelian variety. These twistor lines
  lie in components of $\mathcal{T}_{\omega}$ other than $\mathfrak{H}$.
\end{enumerate}
\end{proposition}

\begin{proof}
(i) The complex structures on $\Lambda_{\RR}$ for which all the tensors
invariant under $G$ are of Hodge type are those $J$ with $J^{2}=-1$ in the
real points of $G\cdot\mathbb{G}_{m}$, since a reductive group containing the
homotheties is the stabiliser of its invariant tensors. Let $\mathfrak{H}$ be
the connected component of this closed set that contains $Y$; it is closed in
$\mathcal{T}$ and consists of the $J=1\otimes1\otimes j_{3}$ with
$j_{3}\in\SL_{2}(\RR)$, $j_{3}^{2}=-1$, in the half plane that contains the
structure of $Y$, that is, of the squares of the members of the Mumford
family. At each of its points the first order condition for $\omega$ to stay
of type $(2,2)$ along $\xi\in H^{1}(T_{Y})$ is $\xi\lrcorner\,\omega=0$, as
in the proof of
\Cref{thm:mumfordrigid}(iii), where no polarisation was used, and by
\Cref{thm:mumfordrigid}(i) the solutions form the line $\CC\kappa$ tangent to
$\mathfrak{H}$. So near each of its points $\mathcal{T}_{\omega}$ is the smooth
curve germ of $\mathfrak{H}$, and $\mathfrak{H}$ is open in
$\mathcal{T}_{\omega}$. Being also closed and connected, it is a connected
component. A connected analytic family on which $\omega$ stays of Hodge type
maps into $\mathcal{T}_{\omega}$, hence into $\mathfrak{H}$ once one of its
points does.

(ii) A twistor line is a nonconstant holomorphic image of $\PP^{1}$ in
$\mathcal{T}$. If $\omega$ stayed of Hodge type along one through a point of
$\mathfrak{H}$, or along any compact connected curve through it, the curve
would lie in $\mathfrak{H}$ by (i), and a holomorphic map from a compact
connected curve to the upper half plane is constant.

(iii) At $\sigma_{1}$ the quaternion algebra $D$ is definite
(\Cref{setup:mumford}), so the first factor of $G(\RR)$ is the group of unit
quaternions, and $V_{\RR}$ restricted to it is a module over
$D\otimes_{F,\sigma_{1}}\RR\cong\mathbb{H}$. The unit quaternions $j$ with
$j^{2}=-1$ form a two-sphere, and the complex structures given by left
multiplication by them on a quaternionic vector space form a twistor line.
Each $j$ lies in $G(\RR)$, and so does the circle $\cos\vartheta+j\sin\vartheta$,
which fixes $\omega$; so $\omega$ is of type $(2,2)$. Inside
$V_{1}\otimes V_{2}\otimes V_{3}$ the space $V_{\RR}$ is the fixed space of a
real structure that is quaternionic on the first two factors and real on the
third, so $V_{\RR}=U\otimes V_{3,\RR}$ with $U$ the real form of
$V_{1}\otimes V_{2}$. On $U$ the form $g=\epsilon_{1}\otimes\epsilon_{2}$ is
symmetric and invariant under the compact group of pairs of unit quaternions,
which acts irreducibly on $U\otimes\CC=V_{1}\otimes V_{2}$; so $g$ is
definite. The complex structures of the Mumford family are $1\otimes h$, with
$h$ a complex structure on $V_{3,\RR}$ for which $\epsilon_{3}(\cdot,h\cdot)$
is definite, and $\psi(x,(1\otimes h)y)=g\otimes\epsilon_{3}(\cdot,h\cdot)$ is
definite. For $J=j\otimes1\otimes1$ the form $g(\cdot,(j\otimes1)\cdot)$ is
alternating, $j\otimes1$ being an isometry of $g$ with square $-1$, and
$\psi(x,Jy)=g(\cdot,(j\otimes1)\cdot)\otimes\epsilon_{3}$ is the product of
nondegenerate alternating forms on spaces of dimensions $4$ and $2$; as the
product of two alternating forms on $\RR^{2}$ is the determinant on
$\RR^{2}\otimes\RR^{2}=M_{2}(\RR)$, of signature $(2,2)$, it has signature
$(4,4)$. At a member whose Mumford--Tate group is $G$, the Hodge
classes of degree two on $V_{\RR}\oplus V_{\RR}$ are the forms $\psi\otimes s$
with $s$ a symmetric $2\times2$ matrix, by \Cref{prop:mumfordproduct}(ii), and
the signature of $\psi(\cdot,J\cdot)\otimes s$ is $(4(p+q),4(p+q))$ when $s$ has
signature $(p,q)$, so none is definite, none is a polarisation, and the member
is not an abelian variety. The members of the Mumford family are not on these
lines, since $\psi(x,Jy)$ is definite at the former and of signature $(4,4)$
at the latter, so by (i) the lines lie in other components.
\end{proof}

Markman's proof for Weil classes on abelian fourfolds moves a sheaf along
twistor lines on which its Chern character stays of Hodge type. By
\Cref{prop:notwistor} no such line passes through the square of a member of
the Mumford family: twistor lines on which $\omega$ is of Hodge type exist,
but they carry non-algebraic tori whose Hodge cocharacter lies in a definite
factor of $D$, and they lie in components of the Hodge locus that the Mumford
family never meets. The Kuga--Satake class of \Cref{thm:mumfordks} behaves
differently. The Kuga--Satake construction applies to every Hodge structure
of K3 type on $(T,q_{\lambda})$, so that class stays of Hodge type over the
whole period domain of $(T,q_{\lambda})$, of dimension seven, and not only
along the curve; among the classes attached to the target, it is the only one
for which deformation methods of this kind are not excluded.

The same space $T$ also makes the square a holomorphic symplectic variety,
and this determines what the mechanism of van Geemen and Rapagnetta would
have to supply here. They prove the Hodge conjecture for the general abelian
fourfold of Weil type with trivial discriminant by pulling back the second
Chern class of a hyperk\"ahler sixfold along a map from the fourfold, which
gives an algebraic class of codimension two that is not a product of divisor
classes \cite{vGR26}. For a projective hyperk\"ahler manifold $M$ write
$\sigma_{M}$ for its symplectic form, $q_{M}$ for its Beauville--Bogomolov
form and $T(M)$ for the orthogonal complement of $\NS(M)_{\QQ}$ in
$H^{2}(M,\QQ)$ for $q_{M}$.

\begin{proposition}[The square as a holomorphic symplectic
variety]\label{prop:hksquare}
Let $c\in C$ be a point at which the Mumford--Tate group of $X=X_{c}$ is
$G\cdot\mathbb{G}_{m}$, let $Y=X\times X$, and let $\Psi\in\bigwedge^{2}V$ be
the $G$-invariant bivector dual to $\psi$. For $x\in T$ put
$\iota(x)=(x\otimes1)\Psi$ in $V\otimes V=H^{1}(X,\QQ)\otimes H^{1}(X,\QQ)$, a
K\"unneth component of $H^{2}(Y,\QQ)$, and let $\iota_{2}\colon S^{2}T\to
H^{4}(Y,\QQ)$ be the map $xy\mapsto\iota(x)\cup\iota(y)$. Let $C_{i}\in
S^{2}T_{i}$ be the Casimir element of the trace form of
$\mathfrak{sl}(V_{i})$, and for $\nu\in F$ let
$c_{\nu}=\sum_{i}\sigma_{i}(\nu)C_{i}$, a rational element of $S^{2}T$; for
$\nu\ne0$ it is the element dual to the form $q_{1/\nu}$.
\begin{enumerate}[label=\textup{(\roman*)},leftmargin=2.6em,itemsep=2pt]
\item $\iota$ is an embedding of Hodge structures $T(-1)\to H^{2}(Y,\QQ)$, its
  image is the only irreducible sub-Hodge structure of $H^{2}(Y,\QQ)$ with
  $h^{2,0}=1$, and the holomorphic $2$-form $\sigma_{0}$ spanning its part of
  type $(2,0)$ is symplectic.
\item $\iota_{2}(c_{\nu})=\rho\,\omega_{a}$ with $a_{0}=3\Tr_{F/\QQ}(\nu)$ and
  $\alpha=4\nu-\Tr_{F/\QQ}(\nu)$, in the notation of \Cref{setup:omegaa},
  where $\rho$ is a nonzero rational number independent of $\nu$. So
  $\iota_{2}(c_{\nu})$ is exceptional if and only if $\nu\notin\QQ$, and every
  rational exceptional class is some $\iota_{2}(c_{\nu})$ plus a rational
  multiple of $\pi_{0}$.
\item Let $M$ be a projective hyperk\"ahler manifold and
  $f\colon Y\dashrightarrow M$ a rational map with $f^{*}\sigma_{M}\ne0$. Then
  $f^{*}$ maps $T(M)$ isomorphically onto $\iota(T)$ and is compatible with
  products of classes in $T(M)$; $f$ is generically finite onto a subvariety
  of dimension eight on which $\sigma_{M}$ is generically nondegenerate;
  $b_{2}(M)\ge10$; and $\dim M\ge10$.
\item Suppose in (iii) that $M=S^{[n]}$ is the Hilbert scheme of $n$ points
  on a projective K3 surface $S$, and
  let $q_{\lambda}$ be the form on $T$ that corresponds to the intersection
  form of $S$ on $T(S)=T(M)$ under $\iota^{-1}f^{*}$
  (\Cref{prop:mumfordrm}(iii)). If $\lambda\notin\QQ$, or if the real
  multiplication of $T(S)$ is induced by algebraic cycles on $S\times S$,
  then the Hodge conjecture holds for $Y$.
\end{enumerate}
\end{proposition}

\begin{proof}
(i) The map $\iota$ is $G$-equivariant because $\Psi$ is invariant, and
injective. Its image lies in $S^{2}V$: for $x$ in the symplectic Lie algebra
of $\psi$ one has $(x\otimes1)\Psi=-(1\otimes x)\Psi$, and the exchange of the
two factors carries $(x\otimes1)\Psi$ to $-(1\otimes x)\Psi=(x\otimes1)\Psi$,
because $\Psi$ is alternating. An element $g$ of $G\cdot\mathbb{G}_{m}$
multiplies $\Psi$ by the multiplier $\chi(g)$ of $\psi$, so $(g\otimes g)\iota(x)=
\chi(g)\,\iota(gxg^{-1})$: the map $\iota$ is equivariant for the adjoint
action on $T$ twisted by $\chi$, which is the action defining $T(-1)$. Both
Hodge structures come from the Hodge cocharacter through
$G\cdot\mathbb{G}_{m}$, so $\iota$ is a morphism of Hodge structures from
$T(-1)$. The sub-Hodge structures of $H^{2}(Y,\QQ)$ are its $G$-stable
subspaces, since $\mathbb{G}_{m}$ acts by scalars. Over $\bar{\QQ}$,
$H^{2}(Y)=\bigwedge^{2}(V\otimes\bar{\QQ}^{2})=S^{2}V\oplus(\bigwedge^{2}V)^{\oplus3}$,
and expanding $V\otimes V=\bigotimes_{k}(S^{2}V_{k}\oplus\bigwedge^{2}V_{k})$
gives $\bigwedge^{2}V$ as the sum of the pieces of highest weights $(0,0,0)$,
$(2,2,0)$, $(2,0,2)$, $(0,2,2)$ (an odd number of factors $\bigwedge^{2}V_{k}$)
and $S^{2}V$ as the sum of those of highest weights $(2,2,2)$, $(2,0,0)$,
$(0,2,0)$, $(0,0,2)$. So the isotypic pieces of $H^{2}(Y)$ have these eight
highest weights and dimensions $3,27,27,27,27,3,3,3$. The Hodge cocharacter
acts through $V_{3}$, on which it has one weight of each type, so $S^{2}V_{3}$
has one dimension of type $(2,0)$ and $\bigwedge^{2}V_{3}$ none; the parts of
type $(2,0)$ of the eight pieces have dimensions $0,0,9,9,9,0,0,1$. The last
three pieces occur once, in $S^{2}V$, and are $\iota(T_{1})$, $\iota(T_{2})$,
$\iota(T_{3})$. So an irreducible summand of highest weight $(2,0,2)$ or
$(0,2,2)$ contributes $3$ to $h^{2,0}$, one of weight $(2,2,2)$ contributes
$9$, one of weight $(0,0,2)$ contributes $1$, and the others contribute
nothing. A rational sub-Hodge structure $W$ is stable under the Galois group,
which permutes the highest weights as it permutes the factors, transitively.
So if $W\otimes\bar{\QQ}$ has a summand of weight $(2,2,0)$, $(2,0,2)$ or
$(0,2,2)$, it has summands of all three and $h^{2,0}(W)\ge6$; if it has one of
weight $(2,2,2)$, then $h^{2,0}(W)\ge9$. If $h^{2,0}(W)=1$, therefore,
$W\otimes\bar{\QQ}$ contains the piece $\iota(T_{3})$ of weight $(0,0,2)$,
which has multiplicity one, and with it $\iota(T_{1})$ and $\iota(T_{2})$; so
$W\supseteq\iota(T)$, and $W=\iota(T)$ if $W$ is irreducible. The same
argument shows that $\iota(T)$ is irreducible: a nonzero sub-Hodge structure
of it contains one of the $\iota(T_{i})$ over $\bar{\QQ}$, hence all three.
The part of type $(2,0)$ of $\iota(T)$ is spanned by $\sigma_{0}=\iota(E_{3})$,
where $E_{3}\in\mathfrak{sl}(V_{3})\otimes\CC$ maps $V_{3}^{0,1}$ onto
$V_{3}^{1,0}$. Write $\Psi=\Psi_{1}\otimes\Psi_{2}\otimes\Psi_{3}$ with
$\Psi_{k}\in\bigwedge^{2}V_{k}$ dual to $\epsilon_{k}$, and
$\Psi_{3}=e_{+}\otimes e_{-}-e_{-}\otimes e_{+}$ with $e_{\pm}$ spanning
$V_{3}^{1,0}$ and $V_{3}^{0,1}$. Then $(E_{3}\otimes1)\Psi_{3}$ is a nonzero
multiple of $e_{+}\otimes e_{+}$, so $\sigma_{0}$ is a nonzero multiple of
$\Psi_{1}\otimes\Psi_{2}\otimes e_{+}\otimes e_{+}$: it lies in
$H^{1,0}(X)\otimes H^{1,0}(X)$ and pairs $H^{1,0}$ of the first factor with
$H^{1,0}$ of the second through the nondegenerate form dual to
$\epsilon_{1}\otimes\epsilon_{2}$ on $V_{1}\otimes V_{2}$. So $\sigma_{0}$ has
rank eight: it is a holomorphic symplectic form on $Y$.

(ii) For $x,y\in T$ the class $\iota(x)\cup\iota(y)$ lies in
$H^{2}(X)\otimes H^{2}(X)$ and comes from the tensor
$(x\otimes1)\Psi\otimes(y\otimes1)\Psi$ in $V^{\otimes4}$ by antisymmetrising
its two slots in the first factor of $Y$ and its two slots in the second.
Read it, as in \Cref{setup:omegaa}, as a map from the first-factor slots to
the second-factor slots. Contraction with $\epsilon_{k}$ turns
$\Psi_{k}\otimes\Psi_{k}$ into $\gamma$ times the identity of
$V_{k}\otimes V_{k}$ and, for $x,y\in T_{k}$, the tensor
$(x\otimes1)\Psi_{k}\otimes(y\otimes1)\Psi_{k}$ into $\gamma\,(x\otimes y)$,
with one and the same constant $\gamma\ne0$, since $x$ and $y$ are
antisymmetric for $\epsilon_{k}$ and the two signs cancel. So $\iota_{2}(C_{1})$
corresponds, up to a nonzero constant $\gamma'$ independent of the factor, to
the restriction to $\bigwedge^{2}V$ of $\mathrm{Cas}_{1}\otimes1\otimes1$,
where $\mathrm{Cas}_{1}=\sum_{c}x_{c}\otimes x^{c}=s_{1}-\frac12$ is the
operator of Step 3 of the proof of \Cref{thm:mumfordpowers}, $s_{1}$ the
exchange of the two copies of $V_{1}$. On $\bigwedge^{2}V$ the exchange
$s_{1}$ is $-1$ on $\QQ\theta$ and $U_{23}$ and $+1$ on $U_{12}$ and $U_{13}$,
so $\iota_{2}(C_{1})=\rho\,(3\pi_{0}-\pi_{12}-\pi_{13}+3\pi_{23})$ with
$\rho=-2\gamma'$, and permuting the factors gives the two cyclic permutations
of this identity. Summing with the coefficients $\sigma_{i}(\nu)$ gives
$\iota_{2}(c_{\nu})=\rho\bigl(3\Tr(\nu)\pi_{0}+\sum\sigma_{k}(4\nu-\Tr(\nu))
\pi_{ij}\bigr)$, the sum over $\{i,j,k\}=\{1,2,3\}$ with $i<j$, which is
$\rho\,\omega_{a}$ with the stated $a_{0}$ and $\alpha$. For rational $\nu\ne0$
both $\iota_{2}(c_{\nu})$ and $\omega_{a}$ are rational and nonzero, so
$\rho\in\QQ$. Since
$\Tr(\nu)\in\QQ$, the element $\alpha$ is rational exactly when $\nu$ is; and
$\Tr(\alpha)=\Tr(\nu)$, so $\nu=(\alpha+\Tr(\alpha))/4$ recovers $\nu$ from
any $\alpha\in F$. Two rational classes with the same $\alpha$ differ by a
multiple of $\pi_{0}$, and $\iota_{2}(c_{\nu/\rho})=\omega_{a}$ because
$c_{\nu}$ is linear in $\nu$; so every rational exceptional class is some
$\iota_{2}(c_{\nu})$ plus a rational multiple of $\pi_{0}$. The Casimir element $C_{i}$ is dual to the trace form
on $T_{i}$, so $c_{\nu}$ is dual to $\sum_{i}\sigma_{i}(\nu)^{-1}
\operatorname{tr}(x_{i}y_{i})=q_{1/\nu}(x,y)$.

(iii) Resolve the indeterminacy of $f$ by a sequence of blow-ups along
smooth centres, $\pi\colon\tilde{Y}\to Y$ and $\tilde{f}\colon\tilde{Y}\to M$,
and put $f^{*}=\pi_{*}\tilde{f}^{*}$, a morphism of Hodge structures. The
kernel of $\pi_{*}$ on $H^{2}(\tilde{Y},\QQ)$ is spanned by the classes of
the exceptional divisors, so $\alpha\mapsto\tilde{f}^{*}\alpha-
\pi^{*}f^{*}\alpha$ is a morphism of Hodge structures from $H^{2}(M,\QQ)$ to a
space of Hodge classes. The Hodge structure $T(M)$ is irreducible: a
sub-Hodge structure $W\subseteq T(M)$ with $W^{2,0}=0$ consists of rational
classes of type $(1,1)$, which lie in $\NS(M)_{\QQ}\cap T(M)=0$; if
$W^{2,0}\ne0$, then $\sigma_{M}\in W_{\CC}$, the orthogonal complement of $W$
in $H^{2}(M,\QQ)$ is a sub-Hodge structure orthogonal to $\sigma_{M}$ and
$\bar{\sigma}_{M}$, hence of type $(1,1)$ and contained in $\NS(M)_{\QQ}$,
and $W\supseteq T(M)$. Since $T(M)$ is not spanned by Hodge classes, that
morphism vanishes on it, and $\tilde{f}^{*}\alpha=\pi^{*}f^{*}\alpha$ for
$\alpha\in T(M)$. Because $\tilde{f}^{*}$ is a ring homomorphism and
$\pi_{*}\pi^{*}$ is the identity, $f^{*}(\alpha\beta)=f^{*}\alpha\cup
f^{*}\beta$ for $\alpha,\beta\in T(M)$. The kernel of $f^{*}$ on $T(M)$ is a
sub-Hodge structure not containing $\sigma_{M}$, hence zero, and the image is
irreducible with $h^{2,0}=1$, hence $\iota(T)$ by (i). So $f^{*}\sigma_{M}$ is
a nonzero multiple of $\sigma_{0}$ and has rank eight at every point of $Y$
where $f$ is defined; the rank of a pulled-back form is at most the dimension
of the image, so $f$ is generically finite onto its image $Z$, of dimension
eight, and $\sigma_{M}$ is generically nondegenerate on $Z$. As $M$ is
projective, $b_{2}(M)=\dim T(M)+\dim\NS(M)_{\QQ}\ge9+1$; this excludes the
generalised Kummer varieties, whose second Betti number is seven
\cite{Bea83}. Finally $\dim M\ge\dim Z=8$; suppose that $\dim M=8$. Then
$f$ is dominant, and for $\alpha\in T(M)_{\CC}$ the relation of Fujiki and
Beauville \cite{Bea83,Fuj87} gives
\[
  \int_{Y}(f^{*}\alpha)^{8}=\int_{\tilde{Y}}\tilde{f}^{*}(\alpha^{8})
  =\deg(f)\int_{M}\alpha^{8}=\deg(f)\,c_{M}\,q_{M}(\alpha)^{4},
  \qquad c_{M}>0 .
\]
Write $f^{*}\alpha=\iota(x)$ with $x=x_{1}+x_{2}+x_{3}\in T_{\CC}$ and
$N_{i}=-\det x_{i}$. The form $x\mapsto q_{M}((f^{*})^{-1}\iota(x))$ is
invariant under $G$, so it is $q_{\lambda}(x,x)=
\sum_{i}\sigma_{i}(\lambda)\operatorname{tr}(x_{i}^{2})=
2\sum_{i}\sigma_{i}(\lambda)N_{i}$ for some $\lambda\in F^{\times}$ by
\Cref{prop:mumfordrm}(iii). On the other side, in a basis $(e_{i})$ of
$V$ one has $\iota(x)=\sum_{i,j}M_{ij}\,e_{i}^{(1)}\wedge e_{j}^{(2)}$, with
$M$ the product of the matrix of $x$ and the constant invertible matrix of
$\Psi$, and the $2$-forms $e_{i}^{(1)}\wedge e_{j}^{(2)}$ commute, so
$\iota(x)^{8}=8!\det(M)\prod_{i}e_{i}^{(1)}\wedge e_{i}^{(2)}$ is $8!\det(x)$
times a fixed volume form, and
$\det(x_{1}\otimes1\otimes1+1\otimes x_{2}\otimes1+1\otimes1\otimes
x_{3})=\Delta(N_{1},N_{2},N_{3})^{2}$ with
$\Delta=N_{1}^{2}+N_{2}^{2}+N_{3}^{2}-2N_{1}N_{2}-2N_{1}N_{3}-2N_{2}N_{3}$;
this is also the product of the eight numbers $\pm r_{1}\pm r_{2}\pm r_{3}$,
where $\pm r_{i}$ are the eigenvalues of $x_{i}$ and $r_{i}^{2}=N_{i}$. The
$N_{i}$ are algebraically independent functions on $T_{\CC}$, so the two
expressions for $\int_{Y}(f^{*}\alpha)^{8}$ would make $\Delta^{2}$ a nonzero
constant multiple of the fourth power of the linear form
$\sum_{i}\sigma_{i}(\lambda)N_{i}$, and then, by unique factorisation,
$\Delta$ a multiple of its square. But $\Delta$ is a nondegenerate
ternary quadratic form, hence irreducible. So $\dim M\ne8$, and since the
dimension of a hyperk\"ahler manifold is even, $\dim M\ge10$.

(iv) By \cite{Bea83} the map $H^{2}(S,\QQ)\to H^{2}(S^{[n]},\QQ)$ induced by
the universal subscheme $Z_{n}\subset S\times S^{[n]}$ is injective, carries
the intersection form to the Beauville--Bogomolov form, and has
$H^{2}(S^{[n]},\QQ)=H^{2}(S,\QQ)\oplus\QQ\delta$ with $\delta$ algebraic, so
$T(M)=T(S)$. Let $\Delta_{T}$ be the K\"unneth component in $T(S)\otimes
T(S)$ of the class of the diagonal of $S$; it is algebraic, being the class
of the diagonal minus $[\mathrm{pt}\times S]$, $[S\times\mathrm{pt}]$ and
$\sum_{r}d_{r}\times d_{r}^{\vee}$ over dual bases of $\NS(S)_{\QQ}$. Its image
under the correspondence $Z_{n}\times Z_{n}$, pulled back along the diagonal
of $M$, is the class $\sum_{a}e_{a}\cup e_{a}^{\vee}$ over dual bases of
$T(S)$ for the intersection form, which is therefore algebraic on $M$, and by
(iii) its image under $f^{*}$ is $\iota_{2}$ of the element of $S^{2}T$ dual
to $q_{\lambda}$, that is $\iota_{2}(c_{1/\lambda})$, an algebraic class on
$Y$. If $\lambda\notin\QQ$ it is exceptional by (ii). If instead the real
multiplication is algebraic, apply it to the first factor of $\Delta_{T}$:
for $b\in F$ the class $\sum_{a}(m_{b}e_{a})\cup e_{a}^{\vee}$ is algebraic,
and its image is $\iota_{2}(c_{b/\lambda})$, which is exceptional for
$b\notin\QQ\lambda$. In both cases \Cref{lem:oneexceptional} makes every
Hodge class in $H^{2}(X)\otimes H^{2}(X)$ algebraic; the Hodge classes in the
other K\"unneth components of $H^{4}(Y,\QQ)$ are products of divisor classes,
as at the end of the proof of \Cref{thm:mumfordks}; and the Hodge classes of
$Y$ form a ring generated in degrees two and four (proof of
\Cref{cor:mumfordequiv}). So every Hodge class of $Y$ is algebraic.
\end{proof}

\begin{remark}[The mechanism of van Geemen and
Rapagnetta]\label{rem:hkroute}
By \Cref{prop:hksquare}(ii) the exceptional classes are, up to rational
multiples of $\pi_{0}$, the classes $\iota_{2}(c_{\nu})$ with $\nu\notin\QQ$,
the images of the elements dual to the forms $q_{\lambda}$, $\lambda=1/\nu$,
on the part $\iota(T)$ of $H^{2}(Y)$ generated by the symplectic class; the
form with $\lambda$ rational gives a product of divisor classes. They are the
analogues for $(Y,\sigma_{0})$ of the dual Beauville--Bogomolov class, whose
pull-back is what \cite{vGR26} uses. By part (iii), a transfer of that
argument needs a projective hyperk\"ahler manifold of dimension at least ten with
transcendental part $T(-1)$, such as $S_{\mu}^{[n]}$ with $n\ge5$ for a
surface $S_{\mu}$ of \Cref{thm:mumfordks}, and a rational map from $Y$ onto a
symplectic subvariety of dimension eight. For $M=S_{\mu}^{[n]}$ the form of
(iv) is $q_{\mu/e^{2}}$ for the $e\in F^{\times}$ with $\iota^{-1}f^{*}=m_{e}$
on $T(S_{\mu})=T$, so $\lambda\notin\QQ$ for every such map as soon as
$\mu\notin\QQ^{\times}F^{\times2}$. No such map is known. It would give an
algebraic cycle on $S_{\mu}\times X\times X$ inducing an isomorphism of
$T(S_{\mu})$ onto $\iota(T)\subset H^{1}(X)\otimes H^{1}(X)$, a correspondence
of the kind that the Kuga--Satake class supplies in the proof of
\Cref{thm:mumfordks}; so this route is the Kuga--Satake route in geometric
form rather than an independent one. By (iii), maps to K3 surfaces and to
hyperk\"ahler eightfolds cannot serve.
\end{remark}

\begin{theorem}[A numerical criterion for the Mumford classes]
\label{thm:mumfordnumerical}
Let $c\in C$, $Y=X_{c}\times X_{c}$, $\omega=\omega_{a}$ a rational
exceptional class with $L_{0}(a)\ne0$, and $E$ a perfect complex on $Y$ with
$\ch(E)=N\omega$, $N\ne0$, and no other component.
\begin{enumerate}[label=\textup{(\roman*)},leftmargin=2.6em,itemsep=2pt]
\item $\dim\Ext^{2}(E,E)\ge119$.
\item If $\dim\Ext^{2}(E,E)=119$, then $\operatorname{ev}_{E}\colon
  HT^{2}(Y)\to\Ext^{2}(E,E)$ is surjective with kernel $\CC\kappa$, and
  $\sigma_{E}$ is injective. Conversely, if $\sigma_{E}$ is injective, then
  $\operatorname{ev}_{E}$ has rank $119$ and kernel $\CC\kappa$, and $E$
  extends to first order along the curve.
\item If some such $E$ at one point $c$ has $\dim\Ext^{2}(E,E)=119$, then
  $\omega_{t}$ is algebraic for every $t\in C$, the Hodge conjecture holds for
  $X_{t}\times X_{t}$ for every $t$, and $B(W\times_{C}W)$ holds.
\item At every CM point $c$ some perfect complex $E$ on $Y$ has
  $\ch(E)=N\omega$ for some $N\ne0$ and no other component.
\end{enumerate}
\end{theorem}

\begin{proof}
(i) and (ii). By the square $\sigma_{E}\circ\operatorname{ev}_{E}=
(\,\cdot\,)\lrcorner\,\ch(E)$ of \Cref{setup:criterion}, $\ker
\operatorname{ev}_{E}\subseteq\Ann_{HT^{2}}(\omega)=\CC\kappa$ by
\Cref{thm:mumfordrigid}(ii), so $\operatorname{ev}_{E}$ has rank at least
$119$, which is (i). If $\dim\Ext^{2}(E,E)=119$, then
$\operatorname{ev}_{E}$ is onto, its kernel has dimension one and lies in
$\CC\kappa$, so equals it; and $\sigma_{E}(\operatorname{ev}_{E}(x))=
N\,x\lrcorner\,\omega$ vanishes only for $x\in\CC\kappa=\ker
\operatorname{ev}_{E}$, so $\sigma_{E}$ is injective on
$\image\operatorname{ev}_{E}=\Ext^{2}(E,E)$. Conversely, if $\sigma_{E}$ is
injective, then $\sigma_{E}(\operatorname{ev}_{E}(\kappa))=N\kappa\lrcorner
\,\omega=0$ forces $\operatorname{ev}_{E}(\kappa)=0$, so the kernel is
$\CC\kappa$ and the rank is $119$; and $\operatorname{ev}_{E}$ restricted to
$H^{1}(T_{Y})$ is the obstruction map \cite{Tod09}, which therefore vanishes
at $\kappa$, so that $E$ extends to first order along the curve.

(iii) By (ii) $\sigma_{E}$ is injective. The Chern character of $E$ is
$N\omega$, a flat section of Hodge type over the whole of $C$, so
\Cref{thm:BF} applies to every infinitesimal deformation of $Y$ inside
$W\times_{C}W\to C$. The proof of \Cref{thm:semiregclosure} then applies word
for word, with the stack of perfect complexes on the fibres of $W\times_{C}
W\to C$ \cite{Lie06} in place of that of $\cA\to\cS$, or the analytic
Kuranishi germ if $\Ext^{<0}(E,E)\ne0$: it is smooth over $C$ at $[E]$, its
image contains a nonempty open subset $U\subseteq C$, and for $t\in U$ a
perfect complex $E_{t}$ on $X_{t}\times X_{t}$ in the same connected
component has $\ch(E_{t})=N\omega_{t}$, so $\omega_{t}$ is algebraic, and by
\Cref{lem:oneexceptional} so are both exceptional classes. Since $U$ is
uncountable, \Cref{cor:mumfordequiv} gives the rest.

(iv) By \Cref{thm:mumfordcm} there is $\xi\in\CH^{2}(Y)_{\QQ}$ with class
$\omega$, and the argument of \Cref{prop:ktheory} gives $E$: the Chern
character is an isomorphism $K_{0}(Y)_{\QQ}\to\CH^{*}(Y)_{\QQ}$
\cite[Ex.~15.2.16(b) and Cor.~18.3.2]{Ful98}, so $\xi=\ch(u)$ for a unique
$u$, and $Nu$ is the class of a perfect complex for a suitable $N$.
\end{proof}

Part (iii) is the Mumford analogue of \Cref{thm:p2numerical}(iii): one object,
at one point, and one number, which is $119$ here and $2n(2n-1)$ for a Weil
family. For a Weil family that number is attained by no object
(\Cref{thm:p2false}), because an object attaining it would have to deform to
non-algebraic tori; the classes here are rigid, so that argument does not
apply (\Cref{rem:p2prime}). The objects of (iv) exist, and the criterion
asks that one of them, or another representative of the same class in
$K_{0}$, attain the bound of (i). The rigidity of \Cref{thm:mumfordrigid} is no obstacle:
semiregularity moves an object along the Hodge locus of its Chern character,
and the curve is exactly that locus, so a semiregular object would propagate
along the curve and nowhere else, which is all that is needed.

\begin{corollary}[Bounded data at CM points]\label{cor:mumfordbounded}
Fix a relatively ample line bundle on $W\times_{C}W\to C$, and let $\omega$ be
a rational exceptional class. The following are equivalent.
\begin{enumerate}[label=\textup{(\alph*)},leftmargin=2.6em,itemsep=2pt]
\item $\omega_{t}$ is algebraic for every $t\in C$; equivalently, by
  \Cref{lem:oneexceptional} and \Cref{cor:mumfordequiv}, the Hodge conjecture
  holds for $X_{t}\times X_{t}$ for every $t\in C$.
\item There are a finite set $\Phi$ of pairs $(\underline{P},\underline{q})$,
  each a tuple of Hilbert polynomials with a tuple of rational numbers of the
  same length, and infinitely many points $c\in C$ at each of which
  $\omega_{c}=\sum_{i}q_{i}[Z_{i}]$ for subschemes $Z_{i}\subset X_{c}\times
  X_{c}$ with Hilbert polynomials $P_{i}$, for some $(\underline{P},
  \underline{q})\in\Phi$.
\end{enumerate}
In particular the classes are algebraic on every member as soon as they are
represented at infinitely many CM points by cycles of bounded degree.
\end{corollary}

\begin{proof}
For a pair $(\underline{P},\underline{q})$ let $\Sigma_{\underline{P},
\underline{q}}\subseteq C$ be the set of points at which $\omega$ is
represented with these data. The proof of \Cref{thm:closedlocus} applies
verbatim to $(W\times_{C}W,C,\omega)$, as in \Cref{lem:spreading}, and shows
that each $\Sigma_{\underline{P},\underline{q}}$ is Zariski closed. If (b)
holds, then, $\Phi$ being finite, one pair of $\Phi$ serves at infinitely
many points, so its $\Sigma$ is an infinite closed subset of the irreducible
curve $C$, hence $C$, and (a) holds. If (a) holds, $C$ is the union of the
countably many $\Sigma_{\underline{P},\underline{q}}$, and $C$ is
uncountable, so one of
them is infinite, hence equal to $C$; then (b) holds with $\Phi$ that one
pair and every point of $C$.
\end{proof}

The proof of \Cref{thm:mumfordcm} represents the classes at every CM point,
through Markman's theorem for abelian fourfolds \cite{Mar25a} and the ring
generated by divisor classes and pull-backs of Hodge classes of $X_{c}$ along
homomorphisms $X_{c}\times X_{c}\to X_{c}$, but with no control of degree;
isogenies between CM members, which move CM points along their Hecke orbits,
do not bound the degree either. By
\Cref{cor:mumfordbounded} such a bound, at infinitely many CM points, is
equivalent to the target; without it the CM points, being countably many,
say nothing about a very general member (\Cref{lem:spreading} needs
uncountably many points).

The next proposition says which endomorphisms those pull-backs have to use.
At a CM point the fourfold is of Weil type, its two Hodge classes of degree
four that are not products of divisor classes form its Weil line, and the
exceptional classes of the square are out of reach of everything built from
that line with the Weil field alone: they need the real multiplication by the
totally real quartic algebra of Rosati-fixed endomorphisms, and all of it.

\begin{proposition}[The Weil structure at a CM point, and what the cycles
must use]\label{prop:cmsource}
Let $c\in C$ be a CM point, $X=X_{c}$, $Y=X\times X$, and let $e_{w}$,
$w\in\{\pm1\}^{3}$, be a weight basis of $V=H^{1}(X,\CC)$ for the Hodge torus
$T$ of $X$, as in the proof of \Cref{thm:mumfordcm}, normalised so that
$\psi(e_{w},e_{-w})=\pm1$, with $e^{(1)}_{w}$, $e^{(2)}_{w}$ its two copies in
$H^{1}(Y,\CC)$. Let $T_{+}$ be the set of the four weights with an even
number of entries $-1$, let $T_{-}=-T_{+}$ be the set of the other four, let
$t_{\pm}=\prod_{w\in T_{\pm}}e_{w}\in H^{4}(X,\CC)$, the factors in a fixed
order, and let $\dagger$ be the Rosati involution of $\theta$ on
$\End^{0}(X)$. For $\varphi\in\End^{0}(X)$ write $\varphi_{w}$ for its
eigenvalue on $e_{w}$.
\begin{enumerate}[label=\textup{(\roman*)},leftmargin=2.6em,itemsep=2pt]
\item $T_{+}$ and $T_{-}$ are the only sets of four weights with sum zero
  that contain no pair of opposite weights, and $\Hdg^{2}(X)\otimes\CC$ is
  the direct sum of the span of the products of divisor classes and of
  $\CC t_{+}\oplus\CC t_{-}$. The elements of $\End^{0}(X)$ with
  $\varphi_{w}$ constant on $T_{+}$ and constant on $T_{-}$ form an imaginary
  quadratic field $k_{c}$; $(X,k_{c},\theta)$ is of $(k_{c},-1,2)$-Weil type,
  and its Weil line is $\HW(X)=(\CC t_{+}\oplus\CC t_{-})\cap H^{4}(X,\QQ)$.
  The fixed algebra $K_{c}^{+}$ of $\dagger$ consists of the elements with
  $\varphi_{w}=\varphi_{-w}$ for every $w$; it is a product of totally real
  fields of total degree four, and $\End^{0}(X)=K_{c}^{+}\otimes_{\QQ}k_{c}$.
\item For $\varphi,\psi\in\End^{0}(X)$ let $f_{\varphi,\psi}^{*}$ be the ring
  homomorphism $H^{*}(X,\QQ)\to H^{*}(Y,\QQ)$ that sends $e\in H^{1}(X,\QQ)$
  to $\mathrm{pr}_{1}^{*}\varphi^{*}e+\mathrm{pr}_{2}^{*}\psi^{*}e$; for
  $\varphi,\psi\in\End(X)$ it is the pull-back along the homomorphism
  $\varphi\circ\mathrm{pr}_{1}+\psi\circ\mathrm{pr}_{2}\colon Y\to X$, and
  every homomorphism $Y\to X$ is of this form. For a $\QQ$-subalgebra
  $R\subseteq\End^{0}(X)$ let $\Sigma_{R}\subseteq\Hdg^{2}(Y)$ be the
  $\QQ$-span of the products of two divisor classes of $Y$ and of the classes
  $f_{\varphi,\psi}^{*}\HW(X)$ with $\varphi,\psi\in R$. Then
  $\Sigma_{k_{c}}=\Sigma_{\QQ}$, the span obtained from the maps
  $(x,y)\mapsto ax+by$ alone, has dimension $110$ and meets the span of
  $\pi_{0},\pi_{12},\pi_{13},\pi_{23}$ in the plane of the classes with
  $a_{1}=a_{2}=a_{3}$; it contains no exceptional class.
\item For a $\QQ$-subalgebra $R\subseteq\End^{0}(X)$ the following are
  equivalent: \textup{(a)} $\Sigma_{R}$ contains an exceptional class;
  \textup{(b)} $\Sigma_{R}=\Hdg^{2}(Y)$, which has dimension $132$;
  \textup{(c)} $R$ separates the weights of $T_{+}$, that is, for any two
  distinct $w,w'\in T_{+}$ some $\varphi\in R$ has
  $\varphi_{w}\ne\varphi_{w'}$; equivalently, $R$ and $k_{c}$ generate
  $\End^{0}(X)$. In particular $\Sigma_{K_{c}^{+}}=\Hdg^{2}(Y)$, while
  $\Sigma_{R}$ contains no exceptional class if $R$ is a proper subalgebra
  of $K_{c}^{+}$ or a subalgebra of $k_{c}$.
\end{enumerate}
\end{proposition}

\begin{proof}
(i) The four weights $d_{1}=(1,1,1)$, $d_{2}=(1,-1,-1)$, $d_{3}=(-1,1,-1)$,
$d_{4}=(-1,-1,1)$ of $T_{+}$ have sum zero, and any three of them are
linearly independent, so every linear relation among them is a multiple of
$d_{1}+d_{2}+d_{3}+d_{4}=0$. A set of four weights containing no opposite
pair contains exactly one of $d_{i}$, $-d_{i}$ for each $i$, say
$\varepsilon_{i}d_{i}$, and has sum zero only if
$\sum_{i}\varepsilon_{i}d_{i}=0$, that is, if all $\varepsilon_{i}$ are
equal: the set is $T_{+}$ or $T_{-}$. The Hodge classes of $X$ of degree
four, tensored with $\CC$, are spanned by the monomials $\prod_{w\in S}e_{w}$
over the sets $S$ of four weights with sum zero (proof of
\Cref{thm:mumfordcm}); the divisor classes are spanned by the four monomials
$e_{w}e_{-w}$ of weight zero and degree two, their products are the six
monomials of the sets that are unions of two opposite pairs, and the two
remaining sets are $T_{+}$ and $T_{-}$.

The torus $T$ is defined over $\QQ$, so the Galois group of $\bar{\QQ}$ acts
linearly on its character lattice and permutes the set of the eight weights
of $V$. A linear map that permutes the weights carries opposite pairs to
opposite pairs and sets with sum zero to sets with sum zero, hence carries
$\{T_{+},T_{-}\}$ to itself. Over $\bar{\QQ}$ the algebra
$\End^{0}(X)\otimes\bar{\QQ}=\End_{T}(V_{\bar{\QQ}})$ consists of the
diagonal matrices $\sum_{w}x_{w}p_{w}$ in the weight basis, $p_{w}$ the
projection to the weight space of $w$, and $\gamma\in\Gal(\bar{\QQ}/\QQ)$
acts by $\gamma(\sum_{w}x_{w}p_{w})=\sum_{w}\gamma(x_{w})p_{\gamma w}$,
because $\gamma$ carries the weight space of $w$ to that of $\gamma w$. As
$\psi$ is $T$-invariant, $\psi(e_{w},e_{w'})=0$ unless $w'=-w$, so the
adjoint of $\sum_{w}x_{w}p_{w}$ for $\psi$, which is the Rosati involution,
is $\sum_{w}x_{-w}p_{w}$. The subalgebra $A_{W}$ of the diagonal matrices
that are constant on $T_{+}$ and on $T_{-}$, and the subalgebra $A^{+}$ of
those with $x_{w}=x_{-w}$ for all $w$, are therefore stable under the Galois
group and under $\dagger$; so $k_{c}=A_{W}\cap\End^{0}(X)$ and
$K_{c}^{+}=A^{+}\cap\End^{0}(X)$ are $\QQ$-forms of $A_{W}$ and $A^{+}$, of
dimensions $2$ and $4$, and $K_{c}^{+}$ is the fixed algebra of $\dagger$.
On $A_{W}$ the involution $\dagger$ exchanges the two scalars. The Rosati
involution is positive: $\Tr(xx^{\dagger})>0$ for $x\ne0$, with $\Tr$ the
trace on $V$ \cite[Section~5.1]{BL04}. If $k_{c}$ were not a field it would
contain the idempotent $p_{+}=\sum_{w\in T_{+}}p_{w}$, and
$\Tr(p_{+}p_{+}^{\dagger})=\Tr(p_{+}p_{-})=0$; so $k_{c}$ is a quadratic
field on which $\dagger$ is the nontrivial automorphism. If $k_{c}$ were
real, $x=\sqrt{d}\in k_{c}$ would satisfy $x^{\dagger}=-x$ and
$\Tr(xx^{\dagger})=-\Tr(x^{2})=-8d<0$. So $k_{c}$ is imaginary quadratic and
$\dagger$ is complex conjugation on it, which is condition (iii) of
\Cref{def:weiltype}. The eigenspaces of $k_{c}\otimes\CC$ on $V$ are
$V_{\pm}=\bigoplus_{w\in T_{\pm}}\CC e_{w}$. The Hodge decomposition of $V$
is the eigenspace decomposition of the Hodge cocharacter, which takes values
in the diagonal torus $T_{\CC}\cdot\mathbb{G}_{m}$, so every $e_{w}$ is of
pure type $(1,0)$ or $(0,1)$, and $t_{+}$ is of type $(p,4-p)$ with
$p=\dim(V_{+}\cap H^{1,0}(X))$. Since $t_{+}$ is a Hodge class, $p=2$: each
eigenvalue of $\sqrt{-d}$ has multiplicity two on $H^{1,0}(X)$, which is
condition (ii) of \Cref{def:weiltype}, and condition (i) is $\dim X=4$. The
Weil line of \Cref{def:weilline} is
$(\bigwedge^{4}V_{+}\oplus\bigwedge^{4}V_{-})\cap H^{4}(X,\QQ)
=(\CC t_{+}\oplus\CC t_{-})\cap H^{4}(X,\QQ)$; it has dimension two by
\Cref{prop:weilline}(i), so it is a rational complement of the products of
divisor classes in $\Hdg^{2}(X)$. The algebra $K_{c}^{+}$ is commutative and
semisimple, being a subalgebra of $\End^{0}(X)$, whose complexification is
commutative, and its trace form $\Tr(xy)=\Tr(xy^{\dagger})$ is positive
definite; so every factor field $E$ of $K_{c}^{+}$ has
$\Tr_{E/\QQ}(x^{2})>0$ for $x\ne0$ and is totally real. Finally the
multiplication map $K_{c}^{+}\otimes_{\QQ}k_{c}\to\End^{0}(X)$ is an
isomorphism, because over $\bar{\QQ}$ it is the isomorphism
$\bar{\QQ}^{\{\text{opposite pairs}\}}\otimes\bar{\QQ}^{\{T_{+},T_{-}\}}
\to\bar{\QQ}^{\{\text{weights}\}}$ induced by the bijection that sends a
weight to the pair and the tetrahedron containing it, a bijection because
the two weights of a pair lie in different tetrahedra, $T_{-}$ being $-T_{+}$.

(ii) Extend the notation: for diagonal $\varphi,\psi\in\End^{0}(X)\otimes\CC$
with entries $\varphi_{w}$, $\psi_{w}$ let $f_{\varphi,\psi}^{*}$ be the ring
homomorphism $H^{*}(X,\CC)\to H^{*}(Y,\CC)$ with
$e_{w}\mapsto\varphi_{w}e^{(1)}_{w}+\psi_{w}e^{(2)}_{w}$; for
$\varphi,\psi\in\End^{0}(X)$ this is the map of the statement, since
$\varphi^{*}e_{w}=\varphi_{w}e_{w}$, and every homomorphism $Y\to X$ is
$\varphi\circ\mathrm{pr}_{1}+\psi\circ\mathrm{pr}_{2}$ for its restrictions
$\varphi,\psi$ to the two factors. Then
\begin{equation}\label{eq:pullbacktet}
  f_{\varphi,\psi}^{*}t_{\pm}
  =\prod_{w\in T_{\pm}}\bigl(\varphi_{w}e^{(1)}_{w}+\psi_{w}e^{(2)}_{w}\bigr)
  =\sum_{S\subseteq T_{\pm}}\pm\Bigl(\prod_{w\in S}\varphi_{w}\Bigr)
  \Bigl(\prod_{w\in T_{\pm}\setminus S}\psi_{w}\Bigr)\,
  e^{(1)}_{S}\,e^{(2)}_{T_{\pm}\setminus S},
\end{equation}
where $e^{(1)}_{S}e^{(2)}_{T\setminus S}$ is the product of the $e^{(1)}_{w}$,
$w\in S$, and the $e^{(2)}_{w}$, $w\in T\setminus S$, taken in the order of
$T$, and the sign is that of the shuffle. We call these $32$ products the
tetrahedron monomials. The Hodge classes of $Y$ of degree four, tensored with
$\CC$, are spanned by the monomials of weight zero in the sixteen classes
$e^{(i)}_{w}$ (proof of \Cref{thm:mumfordcm}). Such a monomial whose four
weights contain an opposite pair is the product of two monomials of weight
zero and degree two, that is, of two divisor classes: there are four of the
form $e^{(1)}_{w}e^{(2)}_{w}e^{(1)}_{-w}e^{(2)}_{-w}$ and $96$ with two
distinct opposite pairs and a choice of copy for each of the four weights.
A weight repeated twice, $2w+w'+w''=0$, forces $w'=w''=-w$, so the remaining
monomials have four distinct weights with no opposite pair, and by (i) they
are the tetrahedron monomials. Hence $\dim\Hdg^{2}(Y)=132$, the products of
divisor classes span the $100$ monomials of the first two kinds, and the
tetrahedron monomials form a basis of the quotient $\Xi$ of
$\Hdg^{2}(Y)\otimes\CC$ by the span of the products of divisor classes. Two
distinct weights of a tetrahedron $T$ differ in exactly two coordinates, and
for $S=\{a,b\}\subseteq T$ the complement $T\setminus S$ has the same set of
differing coordinates, because $a+b$ and the sum of the complement are
opposite. Writing $D(a,b)\subseteq\{1,2,3\}$ for the set of coordinates in
which two weights differ, the six tetrahedron monomials of $T$ with two
factors from each copy fall into three pairs
$\{e^{(1)}_{S}e^{(2)}_{T\setminus S},e^{(1)}_{T\setminus S}e^{(2)}_{S}\}$,
indexed by the two-element set $D(S)=D(T\setminus S)$.

For $\varphi,\psi\in k_{c}\otimes\CC=A_{W}$, with values $\varphi_{\pm}$,
$\psi_{\pm}$ on $T_{\pm}$, \eqref{eq:pullbacktet} reads
$f_{\varphi,\psi}^{*}t_{\pm}=\sum_{k=0}^{4}\varphi_{\pm}^{k}\psi_{\pm}^{4-k}
s_{\pm,k}$, where $s_{\pm,k}$ is the sum, with the signs of
\eqref{eq:pullbacktet}, of the tetrahedron monomials of $T_{\pm}$ with $k$
factors from the first copy. Conversely the classes
$(a\,\mathrm{pr}_{1}+b\,\mathrm{pr}_{2})^{*}t_{\pm}$ for five values of $a/b$
span the $s_{\pm,k}$, a Vandermonde matrix being invertible. Since
$\HW(X)\otimes\CC=\CC t_{+}\oplus\CC t_{-}$, the complexifications of
$\Sigma_{k_{c}}$ and of $\Sigma_{\QQ}$ are both spanned by the products of
divisor classes and the ten classes $s_{\pm,k}$, which are supported on
disjoint nonempty sets of tetrahedron monomials; so
$\Sigma_{k_{c}}=\Sigma_{\QQ}$ has dimension $100+10=110$.

For the four $\pi$ we use the description of \Cref{setup:omegaa}. Fix a
total order on the weights, and for an endomorphism $A$ of
$H^{2}(X,\CC)={\textstyle\bigwedge^{2}}V$ let
\[
  c(A)=\sum_{a<b}A(e_{a}e_{b})\otimes(e_{a}e_{b})^{\vee}
  \in H^{2}(X)\otimes H^{2}(X)
\]
be its class, where $(e_{a}e_{b})^{\vee}$ is the dual basis for the form
${\textstyle\bigwedge^{2}}\psi$. As
${\textstyle\bigwedge^{2}}\psi(e_{a}e_{b},e_{c}e_{d})
=\psi(e_{a},e_{c})\psi(e_{b},e_{d})-\psi(e_{a},e_{d})\psi(e_{b},e_{c})$
vanishes unless $\{c,d\}=\{-a,-b\}$, we have
$(e_{a}e_{b})^{\vee}=\chi(a)\chi(b)\,e_{-a}e_{-b}$ with
$\chi(w)=\psi(e_{w},e_{-w})=\pm1$, the two factors in the order matching
$(a,b)$. Let $s_{t}$, $t=1,2,3$, be the involution of
$V\otimes V=\bigotimes_{u=1}^{3}(V_{u}\otimes V_{u})$ that exchanges the two
factors $V_{t}$. The three $s_{t}$ commute, and $s_{1}s_{2}s_{3}$ is the
exchange of the two slots of $V\otimes V$, which is $-1$ on
${\textstyle\bigwedge^{2}}V$; so each $s_{t}$ preserves
${\textstyle\bigwedge^{2}}V$, on which $s_{1}s_{2}s_{3}=-1$, and
${\textstyle\bigwedge^{2}}V$ is the sum of the four joint eigenspaces with
eigenvalues $(\varepsilon_{1},\varepsilon_{2},\varepsilon_{3})$ of product
$-1$: these are $\CC\theta={\textstyle\bigwedge^{2}}V_{1}\otimes
{\textstyle\bigwedge^{2}}V_{2}\otimes{\textstyle\bigwedge^{2}}V_{3}$ for
$(-1,-1,-1)$ and $U_{ij}=S^{2}V_{i}\otimes S^{2}V_{j}\otimes
{\textstyle\bigwedge^{2}}V_{k}$ for $\varepsilon_{i}=\varepsilon_{j}=1$,
$\varepsilon_{k}=-1$. The projection onto the eigenspace of
$(\varepsilon_{1},\varepsilon_{2},\varepsilon_{3})$ is
$\prod_{t}(1+\varepsilon_{t}s_{t})/2$, which by the relation
$s_{1}s_{2}s_{3}=-1$ equals
$(1+\varepsilon_{1}s_{1}+\varepsilon_{2}s_{2}+\varepsilon_{3}s_{3})/4$. So
\[
  \pi_{0}=c\Bigl(\frac{1-s_{1}-s_{2}-s_{3}}{4}\Bigr),\qquad
  \pi_{12}=c\Bigl(\frac{1+s_{1}+s_{2}-s_{3}}{4}\Bigr),
\]
and $\pi_{13}$, $\pi_{23}$ likewise, with the plus signs at $s_{1},s_{3}$
and at $s_{2},s_{3}$. Let $a^{t}$ be the weight $a$ with its $t$-th
coordinate changed; then $s_{t}(e_{a}e_{b})=e_{a}e_{b}$ if $t\notin D(a,b)$
and $s_{t}(e_{a}e_{b})=e_{a^{t}}e_{b^{t}}$ if $t\in D(a,b)$. Under the
K\"unneth map to $H^{4}(Y,\CC)$ the term of the pair $\{a,b\}$ in $c(s_{t})$
is, up to sign, the monomial $e^{(1)}_{a'}e^{(1)}_{b'}e^{(2)}_{-a}e^{(2)}_{-b}$
with $\{a',b'\}=\{a,b\}$ if $t\notin D(a,b)$ and $\{a',b'\}=\{a^{t},b^{t}\}$
otherwise. If $t\notin D(a,b)$, or $D(a,b)=\{t\}$, in which case $a^{t}=b$,
or $D(a,b)=\{1,2,3\}$, in which case $b=-a$ and the monomial is
$e^{(1)}_{a^{t}}e^{(1)}_{-a^{t}}e^{(2)}_{-a}e^{(2)}_{a}$, the monomial
contains an opposite pair and is a product of two divisor classes. In the
remaining case $D(a,b)=\{t,t'\}$ with $t'\ne t$; then $b=a^{tt'}$, and the
four weights $a^{t}$, $b^{t}=a^{t'}$, $-a$, $-b=a^{t''}$, where
$\{t,t',t''\}=\{1,2,3\}$, are pairwise non-opposite, since any two of them
differ in exactly two coordinates, and they have sum zero, because
$a^{t}+b^{t}$ and $a+b$ agree outside the $t$-th coordinate, where both
vanish. So
\[
  m_{ab}=\chi(a)\chi(b)\,e^{(1)}_{a^{t}}e^{(1)}_{b^{t}}e^{(2)}_{-a}e^{(2)}_{-b},
  \qquad t=\min D(a,b),
\]
is a tetrahedron monomial with $D(a^{t},b^{t})=D(-a,-b)=D(a,b)$, and since
$s_{t'}(e_{a}e_{b})=e_{a^{t'}}e_{b^{t'}}=e_{b^{t}}e_{a^{t}}=-e_{a^{t}}e_{b^{t}}$,
the term of the pair in $c(s_{t'})$ is $-m_{ab}$, while its term in
$c(s_{t''})$ is a product of divisor classes. For each two-element set $D$
put $u_{D}=\sum_{D(a,b)=D}m_{ab}$, a sum over the four pairs with difference
set $D$: two of these pairs lie in each tetrahedron and $a^{t}$ lies in the
other one, so $u_{D}$ has two monomials in each tetrahedron, all with
coefficients $\pm1$, and the three $u_{D}$ are supported on disjoint sets of
monomials; the two monomials of $u_{D}$ in $T$ are those of the pair of
$T$ indexed by $D$. In $\Xi$ this gives $c(1)\equiv0$,
$c(s_{1})\equiv u_{12}+u_{13}$, $c(s_{2})\equiv-u_{12}+u_{23}$,
$c(s_{3})\equiv-u_{13}-u_{23}$, hence
\begin{equation}\label{eq:pimodulo}
  \pi_{0}\equiv0,\quad
  \pi_{12}\equiv\tfrac12(u_{13}+u_{23}),\quad
  \pi_{13}\equiv\tfrac12(u_{12}-u_{23}),\quad
  \pi_{23}\equiv-\tfrac12(u_{12}+u_{13}),
\end{equation}
and $\omega_{a}\equiv\frac12\bigl((a_{2}-a_{3})u_{12}+(a_{1}-a_{3})u_{13}
+(a_{1}-a_{2})u_{23}\bigr)$ in $\Xi$: on the monomials of the pair of
$T_{\pm}$ indexed by $\{1,2\}$, $\{1,3\}$, $\{2,3\}$ the class $\omega_{a}$
has coefficients of absolute values $|a_{2}-a_{3}|/2$, $|a_{1}-a_{3}|/2$,
$|a_{1}-a_{2}|/2$.

Now suppose that $\omega_{a}\in\Sigma_{\QQ}\otimes\CC$. Its image in $\Xi$
is a combination of the $s_{\pm,k}$, and since every monomial of
$\omega_{a}$ has two factors from each copy, only $s_{+,2}$ and $s_{-,2}$
occur: $\omega_{a}\equiv x\,s_{+,2}+y\,s_{-,2}$. On the six monomials of
$T_{+}$ the right side has coefficients of absolute value $|x|$. So the
three differences $a_{2}-a_{3}$, $a_{1}-a_{3}$, $a_{1}-a_{2}$ have a common
absolute value $r$, and as $a_{1}-a_{3}=(a_{1}-a_{2})+(a_{2}-a_{3})$,
$r\in\{0,2r\}$; hence $r=0$ and $a_{1}=a_{2}=a_{3}$. Conversely the classes
with $a_{1}=a_{2}=a_{3}$ are the combinations of $\pi_{0}$ and
$\pi_{12}+\pi_{13}+\pi_{23}$, which are invariant under $\Sp(V,\psi)$ and so
are combinations of the products of divisor classes $\psi_{11}\psi_{22}$ and
$\psi_{12}^{2}$ (proof of \Cref{prop:mumfordrm}(ii)); they lie in
$\Sigma_{\QQ}$. An exceptional class
has for $(a_{1},a_{2},a_{3})$ the three distinct conjugates of an element of
$F\setminus\QQ$ (\Cref{lem:conjugates}), so it does not lie in
$\Sigma_{\QQ}$.

(iii) We first show that (c) is equivalent to $R$ and $k_{c}$ generating
$\End^{0}(X)$, and that $R$ separates the weights of $T_{+}$ if and only if
it separates those of $T_{-}$. For $\varphi\in\End^{0}(X)$ the eigenvalue on
$e_{-w}$ is $\varphi_{-w}=(\varphi^{\dagger})_{w}$; since $\End^{0}(X)$ is a
product of CM fields on which $\dagger$ is complex conjugation, and every
embedding of a CM field into $\CC$ commutes with complex conjugation,
$\varphi_{-w}=\bbar{\varphi_{w}}$. Hence $\varphi_{w}=\varphi_{w'}$ if and
only if $\varphi_{-w}=\varphi_{-w'}$, and $R$ separates $T_{+}$ if and only
if it separates $T_{-}$. The elements of $k_{c}$ are constant on each
tetrahedron and separate the two, and a subalgebra of
$\bar{\QQ}^{\{\text{weights}\}}$ containing the constants is the algebra of
functions constant on the blocks of the partition it induces; so $R$
separates $T_{+}$ exactly when $R\cdot k_{c}\otimes\bar{\QQ}$ separates all
eight weights, that is, when $R\cdot k_{c}=\End^{0}(X)$.

(c) implies (b). If $R$ separates $T_{+}$, then $R\otimes\CC$ maps onto
$\CC^{T_{+}}$ under $\varphi\mapsto(\varphi_{w})_{w\in T_{+}}$, so the
entries $\varphi_{w}$, $\psi_{w}$, $w\in T_{+}$, are eight independent
coordinates on the image of $(R\otimes\CC)^{2}$. The sixteen coefficients in
\eqref{eq:pullbacktet} are distinct monomials in these coordinates, hence
linearly independent functions of $(\varphi,\psi)\in(R\otimes\CC)^{2}$; so
no nonzero linear form on the span of the sixteen tetrahedron monomials of
$T_{+}$ vanishes on all $f_{\varphi,\psi}^{*}t_{+}$, and these classes span
that space. The same holds for $T_{-}$, and with the $100$ products of
divisor classes this is all of $\Hdg^{2}(Y)\otimes\CC$. The map
$(\varphi,\psi)\mapsto f_{\varphi,\psi}^{*}t_{\pm}$ is polynomial and
$R\times R$ is Zariski dense in $(R\otimes\CC)^{2}$, so a linear form
vanishing on its values at rational pairs vanishes on all values, and
finitely many pairs $\varphi,\psi\in R$ give the same span. The $\QQ$-span of
the rational classes $f_{\varphi,\psi}^{*}\HW(X)$, $\varphi,\psi\in R$, and
of the products of divisor classes is then a $\QQ$-subspace of
$\Hdg^{2}(Y)$ whose complexification is the whole space, hence is
$\Hdg^{2}(Y)$.

(b) implies (a), as $\Hdg^{2}(Y)$ contains the exceptional classes.

(a) implies (c). Suppose that two distinct weights $u,u'\in T_{+}$ have
$\varphi_{u}=\varphi_{u'}$ for every $\varphi\in R$, and let $v,v'$ be the
other two weights of $T_{+}$. In \eqref{eq:pullbacktet} the sets
$S=\{u,v\}$ and $S'=\{u',v\}$ then have the same coefficient for every
$\varphi,\psi\in R$, up to the signs of the shuffles, which do not depend on
$\varphi,\psi$. Hence every class in $\Sigma_{R}\otimes\CC$ has coefficients
of equal absolute value on the tetrahedron monomials $e^{(1)}_{S}e^{(2)}_{T_{+}\setminus S}$ and $e^{(1)}_{S'}e^{(2)}_{T_{+}\setminus S'}$, the products
of divisor classes involving no tetrahedron monomial. The pairs
$\{S,T_{+}\setminus S\}$ and $\{S',T_{+}\setminus S'\}$ are different
partitions of $T_{+}$, so $D(S)\ne D(S')$, and for an exceptional class
$\omega_{a}$ in $\Sigma_{R}$ two of the three numbers $|a_{2}-a_{3}|$,
$|a_{1}-a_{3}|$, $|a_{1}-a_{2}|$ would be equal. Each of the six resulting
equations, $a_{2}-a_{3}=\pm(a_{1}-a_{3})$ and the four others, is a linear
relation $c_{1}a_{1}+c_{2}a_{2}+c_{3}a_{3}=0$ with rational coefficients not
all equal, which \Cref{lem:conjugates} forbids for the conjugates of an
element of $F\setminus\QQ$. So (a) fails, and this proves the equivalence.

For the last assertions: $K_{c}^{+}\otimes\bar{\QQ}$ consists of all
functions on the opposite pairs, and $T_{+}$ meets every pair once, so
$K_{c}^{+}$ separates $T_{+}$; a proper subalgebra $R$ of $K_{c}^{+}$ has
$R\otimes\bar{\QQ}$ a proper subalgebra of the functions on the four pairs,
which identifies two pairs and hence two weights of $T_{+}$; and a
subalgebra of $k_{c}$ is constant on $T_{+}$.
\end{proof}

\Cref{fig:cmcube} draws the weights of part (i), and \Cref{fig:cmsource}
the Hodge classes of parts (ii) and (iii).

\begin{figure}[tbp]
\centering
  \includegraphics[width=0.96\textwidth]{fig_cmcube}
\caption{\Cref{prop:cmsource}(i). The eight weights $(\pm1,\pm1,\pm1)$ of the
Hodge torus at a CM point are the vertices of a cube. The four diagonals are
the pairs of opposite weights, paired by the polarisation; the two inscribed
tetrahedra $T_{+}$ (an even number of entries $-1$) and $T_{-}=-T_{+}$ are
the only zero-sum four-sets without an opposite pair, and $t_{\pm}=\prod_{w\in
T_{\pm}}e_{w}$ span the Weil line. The Galois group acts by linear maps
permuting the vertices, so it preserves the diagonals and the pair of
tetrahedra: the functions constant on the tetrahedra form the imaginary
quadratic Weil field $k_{c}$, those constant on the diagonals the totally real
quartic algebra $K_{c}^{+}$ of Rosati-fixed endomorphisms, and
$\End^{0}(X_{c})=K_{c}^{+}\otimes k_{c}$ because a vertex is determined by
its diagonal and its tetrahedron.}
\label{fig:cmcube}
\end{figure}

\begin{figure}[tbp]
\centering
  \includegraphics[width=\textwidth]{fig_cmsource}
\caption{\Cref{prop:cmsource}(ii) and (iii). Top: the $132$ Hodge classes of
degree four on the square at a CM point, one cell per monomial of weight
zero. The $100$ products of divisor classes are the cells with an opposite
pair; the $32$ tetrahedron monomials $e^{(1)}_{S}e^{(2)}_{T\setminus S}$ are
arranged by tetrahedron and by the number $|S|$ of factors from the first
copy. Pull-backs of the Weil line along homomorphisms with components in
$k_{c}$ add only the ten symmetric sums $s_{\pm,k}$, one per column, giving
$110$; components separating the four weights of a tetrahedron, as those of
$K_{c}^{+}$ do, separate the monomials and give all $132$. Bottom: on the six
mixed monomials of one tetrahedron, grouped in the three pairs by the
coordinates in which the two weights of $S$ differ, the symmetric sum
$s_{+,2}$ has coefficients $\pm1$ throughout, while an exceptional class
$\omega_{a}$ has coefficients of absolute value $|a_{2}-a_{3}|/2$,
$|a_{1}-a_{3}|/2$, $|a_{1}-a_{2}|/2$, drawn for $\pi_{12}-\pi_{13}$; two
equal absolute values would be a linear relation among the conjugates of an
irrational element of $F$.}
\label{fig:cmsource}
\end{figure}

The CM points fall into two kinds, according to whether the fourfold is
simple; at the points of the second kind the real multiplication comes from
the fixed cubic field $F$.

\begin{lemma}[Split CM points]\label{lem:splitcm}
Let $c\in C$ be a CM point, and let $L\subset D$ be the quadratic extension
of $F$ whose norm one torus $\Res_{F/\QQ}\ker(\Nm_{L/F})$ is the Hodge group
of $X_{c}$ (proof of \Cref{thm:mumfordcm}). Then $L$ is a CM field, and the
following are equivalent: \textup{(a)} $X_{c}$ is not simple; \textup{(b)}
$K_{c}^{+}$ is not a field; \textup{(c)} $L=Fk$ for an imaginary quadratic
field $k\subseteq L$. In that case $k_{c}\cong k$, $K_{c}^{+}\cong\QQ\times F$,
and $X_{c}$ is isogenous to $E\times B$, where $E$ is an elliptic curve with
complex multiplication by $k$ and $B$ is an abelian threefold of CM type on
which the factor $F$ of $K_{c}^{+}$ acts, a real multiplication by $F$; the
Weil line $\HW(X_{c})$ lies in $H^{1}(E)\otimes H^{3}(B)$. At every other CM
point $K_{c}^{+}$ is a totally real quartic field. Points of the second kind
exist for infinitely many $k$: for every imaginary quadratic field $k$ such
that $k\otimes_{\QQ}F_{v}$ is a field at each finite place $v$ of $F$ at
which $D$ is ramified.
\end{lemma}

\begin{proof}
Write $T$ for the Hodge group. By \Cref{prop:cmsource}(i), $\End^{0}(X_{c})
=K_{c}^{+}\otimes k_{c}$ is a product of CM fields on which $\dagger$ is
complex conjugation, and $T$ lies in the centraliser of $\End^{0}(X_{c})$,
which over $\bar{\QQ}$ is the diagonal torus, and in $\Sp(V,\psi)$; so
$T(\RR)$ lies in the compact group of the $x\in(\End^{0}(X_{c})\otimes\RR)
^{\times}$ with $xx^{\dagger}=1$. If $L\otimes_{F,\sigma}\RR$ were
$\RR\times\RR$ at a real place $\sigma$, the factor of $T(\RR)$ at $\sigma$
would be the norm one group of $\RR\times\RR$ over $\RR$, which is not
compact. Hence $L\otimes_{F,\sigma}\RR=\CC$ at each real place, and $L$ is a
CM field.

Over $\bar{\QQ}$, $T=\prod_{\sigma}T_{\sigma}$ with $T_{\sigma}
=\ker(\Nm_{L/F})\otimes_{F,\sigma}\bar{\QQ}$ the diagonal torus of
$\SL_{2}(\bar{\QQ})=(D\otimes_{F,\sigma}\bar{\QQ})^{1}$ for the decomposition
$L\otimes_{F,\sigma}\bar{\QQ}=\bar{\QQ}\times\bar{\QQ}$ by the two embeddings
of $L$ over $\sigma$, one of them chosen and called $\tau_{\sigma}$, the
other $\bar{\tau}_{\sigma}$. The factor $V_{\sigma}$ of
$V=\bigotimes_{\sigma}V_{\sigma}$ (\Cref{setup:mumford}) has the weights
$\tau_{\sigma}$ and $\bar{\tau}_{\sigma}=\tau_{\sigma}^{-1}$ on $T_{\sigma}$,
with eigenvectors $v_{\sigma,\pm}$, and $e_{w}=\bigotimes_{\sigma}
v_{\sigma,w_{\sigma}}$. An element $\gamma$ of the Galois group carries
$V_{\sigma}$ to $V_{\gamma\sigma}$ and $v_{\sigma,+}$ to
$v_{\gamma\sigma,\varepsilon_{\gamma}(\sigma)}$, where
$\varepsilon_{\gamma}(\sigma)=1$ if $\gamma\tau_{\sigma}=\tau_{\gamma\sigma}$
and $-1$ if $\gamma\tau_{\sigma}=\bar{\tau}_{\gamma\sigma}$; so it acts on
the weights by $(\gamma w)_{\gamma\sigma}=\varepsilon_{\gamma}(\sigma)
w_{\sigma}$.

If $L=Fk$, then $\operatorname{Hom}(L,\bar{\QQ})=\operatorname{Hom}(F,\bar{\QQ})
\times\operatorname{Hom}(k,\bar{\QQ})$ with the Galois group acting on both
factors, and with $\tau_{\sigma}=\sigma\otimes\iota_{0}$ for a fixed
embedding $\iota_{0}$ of $k$ the sign $\varepsilon_{\gamma}(\sigma)
=\varepsilon_{\gamma}$ does not depend on $\sigma$: it is $1$ if $\gamma$
fixes $k$ and $-1$ otherwise. Then $\gamma$ permutes the coordinates of a
weight and changes all its signs when it moves $k$. It preserves each
tetrahedron if it fixes $k$ and exchanges the two otherwise, so
$k_{c}\otimes\bar{\QQ}=\bar{\QQ}^{\{T_{+},T_{-}\}}$ is the algebra of the
Galois set $\operatorname{Hom}(k,\bar{\QQ})$, and $k_{c}\cong k$. It carries
the pair $\{\pm(1,1,1)\}$ to itself, so the projection $p_{E}$ onto
$\CC e_{(1,1,1)}\oplus\CC e_{(-1,-1,-1)}$ lies in $\End^{0}(X_{c})$, and by
Poincar\'e reducibility $X_{c}$ is isogenous to $E\times B$ with
$H^{1}(E,\CC)$ the image of $p_{E}$, of dimension two, so $E$ is an
elliptic curve, on whose $H^{1}$ the field $k_{c}$ acts by two conjugate
scalars: $E$ has complex multiplication by $k$. The threefold $B$ has
$\End^{0}(B)\otimes\bar{\QQ}$ containing the diagonal algebra of the six
remaining weights, of dimension $6=2\dim B$, so $B$ is of CM type. Each of
the six remaining weights has exactly one coordinate different from the
other two, and $\gamma$ carries the opposite pair of the weight with
distinguished coordinate $t$ to the pair with distinguished coordinate
$\gamma t$; so the Galois set of the four opposite pairs is the fixed point
$\{\pm(1,1,1)\}$ together with a copy of $\operatorname{Hom}(F,\bar{\QQ})$,
and $K_{c}^{+}\cong\QQ\times F$, the factor $\QQ$ acting on $H^{1}(E)$ and
the factor $F$ on $H^{1}(B)$. Finally $T_{+}$ consists of $(1,1,1)$ and
three weights of $B$, and $T_{-}$ of $(-1,-1,-1)$ and three weights of $B$,
so $t_{\pm}\in H^{1}(E)\otimes H^{3}(B)$ and $\HW(X_{c})$ lies there. This
proves that (c) implies (a) and (b), with the asserted structure.

$X_{c}$ is simple exactly when $\End^{0}(X_{c})=K_{c}^{+}\otimes k_{c}$ is a
field, and since every factor field $E'$ of $K_{c}^{+}$ is totally real,
$E'\otimes k_{c}$ is a field; so (a) and (b) are equivalent. Suppose (b).
Then the Galois group does not act transitively on the four opposite pairs,
whose algebra is $K_{c}^{+}$. It acts transitively on the three
coordinates, as on the embeddings of $F$, so some $\gamma$ permutes them
cyclically. For each coordinate $t$ the four pairs fall into two sets of
two, according to the value of the coordinate $t$ on the representatives in
$T_{+}$, and this assigns the three coordinates bijectively to the three
partitions of the four pairs into two sets of two, compatibly with the
Galois action, which preserves or exchanges $T_{+}$ and $T_{-}$. So the
permutation of the pairs induced by $\gamma$ has order divisible by three,
hence equal to three: it is a cycle of length three and fixes exactly one
pair. The orbits of the Galois group on the pairs are therefore either one
orbit or two, of sizes one and three, and by hypothesis some pair
$\{w,-w\}$ is fixed: $\gamma w=\delta_{\gamma}w$ with
$\delta_{\gamma}=\pm1$ for every $\gamma$. Replace the choice of
$\tau_{\sigma}$ by $\tau'_{\sigma}=\tau_{\sigma}$ if $w_{\sigma}=1$ and
$\tau'_{\sigma}=\bar{\tau}_{\sigma}$ if $w_{\sigma}=-1$; for the new choice
$w$ becomes $(1,1,1)$, and the new signs are
$\varepsilon'_{\gamma}(\sigma)=\varepsilon_{\gamma}(\sigma)w_{\sigma}
w_{\gamma\sigma}=\delta_{\gamma}$ for all $\sigma$. So $\gamma$ carries the
set $\Phi=\{\tau'_{\sigma}\}$ of three embeddings of $L$ to $\Phi$ if
$\delta_{\gamma}=1$ and to $\bar{\Phi}$ if $\delta_{\gamma}=-1$. The map
$\gamma\mapsto\delta_{\gamma}$ is a character of the Galois group; it is
nontrivial, for otherwise $\Phi$ would be a Galois stable set of three of the
six embeddings and $L$ would not be a field; and complex conjugation carries
every embedding $\tau$ of the CM field $L$ to $\bar{\tau}$, hence $\Phi$ to
$\bar{\Phi}$, so the quadratic field $k$ cut out by the character is
imaginary. Let $\iota_{0}$, $\bar{\iota}_{0}$ be the embeddings of $k$, with
$\iota_{0}$ fixed by the kernel of the character. Then $\tau\mapsto(\tau|_{F},
\iota_{0}\text{ if }\tau\in\Phi\text{ and }\bar{\iota}_{0}\text{ otherwise})$
is a Galois equivariant bijection $\operatorname{Hom}(L,\bar{\QQ})\to
\operatorname{Hom}(F,\bar{\QQ})\times\operatorname{Hom}(k,\bar{\QQ})$, so
$L\cong F\otimes_{\QQ}k$, that is, $L=Fk$ for a copy of $k$ in $L$, which is
(c). At a CM point where (c) fails, $K_{c}^{+}$ is therefore a field, quartic
and totally real by \Cref{prop:cmsource}(i).

Existence. Let $k$ be imaginary quadratic. By the theorem of Albert, Brauer,
Hasse and Noether \cite[\S1.5]{PR94}, the CM field $Fk$ splits $D$, and so
embeds in $D$ over $F$, exactly when $Fk\otimes_{F}F_{v}$ is a field at
every place $v$ of $F$ at which $D$ is ramified, a quadratic extension of a
local field splitting the local quaternion division algebra. At the two real
places where $D$ is definite this holds because $k$ is imaginary, and at the
finite ramified places it is the hypothesis. Given the finitely many
residue characteristics $p$ of these places, choose for each a quadratic
extension $k_{p}$ of $\QQ_{p}$ not contained in $F_{v}$ for any $v$ above
$p$: this is possible since the completions $F_{v}$ above $p$ have degrees
over $\QQ_{p}$ adding up to three, so at most one of them contains a
quadratic extension of $\QQ_{p}$, and then exactly one, while $\QQ_{p}$ has
at least three. Infinitely many imaginary quadratic fields $k$ have the
completions $k_{p}$ at these $p$, and for them
$F_{v}\otimes_{\QQ_{p}}k_{p}$ is a field. For such $k$ let
$T_{k}=\Res_{F/\QQ}\ker(\Nm_{Fk/F})\subset G$, a maximal torus with
$T_{k}(\RR)$ compact. A compact subgroup of $G(\RR)$ fixes a point $x$ of the
symmetric space of $G(\RR)$, and the stabiliser of $x$ is the centraliser of
$h_{x}(\mathbb{S})$; so $h_{x}(\mathbb{S})$ centralises $T_{k}(\RR)$, which
is Zariski dense in $T_{k}$, and lies in $T_{k}(\RR)$. The Hodge group of
the fibre $X_{c}$ over the image $c$ of $x$ is the smallest $\QQ$-subgroup
of $G$ containing $h_{x}(U(1))$, so it lies in $T_{k}$; it has rank three
(proof of \Cref{thm:mumfordcm}), so it equals $T_{k}$, and $c$ is a CM
point with $L=Fk$.
\end{proof}

\begin{remark}[What the degree bound has to control]\label{rem:cmsource}
\leavevmode
\begin{enumerate}[label=\textup{(\arabic*)},leftmargin=2.6em,itemsep=3pt]
\item Markman's theorem \cite{Mar25a}, which the proof of
  \Cref{thm:mumfordcm} invokes, produces the cycles on $X_{c}$ through its
  Weil structure for $k_{c}$. \Cref{prop:cmsource} says that this structure
  is blind to the exceptional classes of the square: everything obtained
  from the Weil line and from divisors by pull-back along homomorphisms with
  components in $k_{c}$, in particular along the maps $(x,y)\mapsto ax+by$,
  and by products, spans $110$ of the $132$ Hodge classes and misses every
  exceptional class; and a representation of an exceptional class by
  pull-backs of the Weil line and products of divisors must use
  homomorphisms whose components, with $k_{c}$, generate the whole
  endomorphism algebra. At a CM point the exceptional classes are thus a
  statement about the real multiplication $K_{c}^{+}$, as they are at a
  general point about the real multiplication by $F$
  (\Cref{prop:mumfordrm}).
\item For $\varphi,\psi\in\End(X_{c})$ and a cycle $Z$ of codimension two on
  $X_{c}$, the degree $\int_{Y}f_{\varphi,\psi}^{*}[Z]\cup\theta_{Y}^{6}$
  is, by \eqref{eq:pullbacktet} applied to the monomials of $[Z]$, a
  polynomial of degree four in the eigenvalues $\varphi_{w}$, $\psi_{w}$,
  which are algebraic integers. For cycles of this shape the bound of
  \Cref{cor:mumfordbounded} therefore asks, at infinitely many CM points,
  for endomorphisms with bounded eigenvalues whose components generate
  $\End^{0}(X_{c})$ over $k_{c}$. Only finitely many algebraic integers of
  degree at most eight have bounded conjugates, so the same finitely many
  minimal polynomials would have to occur at infinitely many points. At
  the split points of \Cref{lem:splitcm} the algebra $K_{c}^{+}=\QQ\times F$
  does not depend on the point and the real multiplication is by $F$ on the
  threefold factor; the representation of the proof of \Cref{thm:mumfordcm}
  then uses, besides elements of $F$, the isogeny $X_{c}\to E\times B$, a
  cycle representing the Weil line of $E\times B$ in
  $H^{1}(E)\otimes H^{3}(B)$, and divisor classes, and these three are what
  varies with the point.
\item At a split point the Weil line is a Hodge class in $H^{1}(E)\otimes
  H^{3}(B)$, a correspondence between an elliptic curve with complex
  multiplication by $k$ and a CM threefold with real multiplication by $F$,
  and its algebraicity is again Markman's theorem for the fourfold
  $E\times B$. A construction of that correspondence with a degree bounded
  independently of $k$, together with a bound on the isogeny and on the
  divisor classes, would give the target through \Cref{cor:mumfordbounded}.
  We do not have one.
\end{enumerate}
\end{remark}

\begin{remark}[Routes to the Mumford target]\label{rem:bypassaudit}
We list the routes to the Mumford target that avoid the Lefschetz standard
conjecture, with what is known about each; \Cref{fig:bypass} summarises
them.
\begin{enumerate}[label=\textup{(\arabic*)},leftmargin=2.6em,itemsep=2pt]
\item \emph{Degeneration.} Closed off: the base is compact and the Hodge
  locus is the base (\Cref{cor:nodegeneration}), so no family on which
  $\omega$ stays of Hodge type approaches the boundary, and there is no limit
  mixed Hodge structure to specialise to.
\item \emph{Deformation into a larger Hodge locus}, on which $\omega$ would
  become a product of divisor classes or a Weil class, hence algebraic by a
  known theorem, and could then be carried back. Closed off by
  \Cref{thm:mumfordrigid}(iii), already to first order.
\item \emph{Hochschild deformations}: twisting $D^{b}(Y)$ by a $B$-field or a
  Poisson bivector, as in the gerby and noncommutative deformations used for
  Weil classes. Closed off: an object with $\ch(E)=N\omega$ extends along no
  direction of $HH^{2}(Y)$ except that of the curve
  (\Cref{cor:mumfordsplit}(i)).
\item \emph{Constructions from line bundles}: cones, tensor products, duals,
  pull-backs and push-forwards along homomorphisms, and Fourier--Mukai
  functors whose kernels are built in the same way. Closed off at every point
  of the curve: everything so built has a Chern character invariant under the
  Lefschetz group, and the exceptional classes are not
  (\Cref{thm:lefschetzclosure}).
\item \emph{Twistor and hyperholomorphic deformation}, the mechanism of
  Markman's proof for Weil classes: move a sheaf along twistor lines on which
  its Chern character stays of Hodge type. Closed off on the square: no
  twistor line through a member of the Mumford family keeps $\omega$ of Hodge
  type, and the twistor lines that do keep it of Hodge type carry
  non-algebraic tori and lie in other components of the Hodge locus
  (\Cref{prop:notwistor}).
\item \emph{Specialisation.} The class is algebraic at every CM point
  (\Cref{thm:mumfordcm}), on every power of $X_{c}$ \cite[Theorem~1.3]{Li26};
  this suffices exactly when the degrees are bounded
  (\Cref{cor:mumfordbounded}), and no bound is known. The cycles that
  represent the class at a CM point must use the real multiplication by the
  Rosati-fixed algebra $K_{c}^{+}$, all of it; the Weil structure of the
  fourfold, with which Markman's theorem produces its cycles, reaches $110$
  of the $132$ Hodge classes and none of the exceptional ones
  (\Cref{prop:cmsource}, \Cref{rem:cmsource}).
\item \emph{Reduction modulo $p$.} At every closed point of every reduction
  of the curve the class is a $\QQ_{\ell}$-combination of algebraic cycles
  (\Cref{prop:modp}), by Li's theorem \cite[Theorem~1.1]{Li26}; this
  suffices, with no lifting step, exactly when the Hilbert data of the cycles
  are bounded on a Zariski dense set of closed points of the arithmetic model
  (\Cref{thm:modpdensity}), and no bound is known.
\item \emph{Semiregularity.} Reduced to one dimension count at one CM point
  (\Cref{thm:mumfordnumerical}); an object attaining it has large odd
  self-extensions (\Cref{thm:mumfordprofile}), is not built from line
  bundles (\Cref{prop:mumforddivisor}) and splits only in a constrained way
  (\Cref{cor:mumfordsplit}); no such object is known.
\item \emph{Kuga--Satake.} Reduced to the Kuga--Satake class of one K3 surface
  of Picard number thirteen (\Cref{thm:mumfordks}), which the known cases of
  the Kuga--Satake Hodge conjecture do not reach (\Cref{rem:mumfordks}). Its
  geometric form is the mechanism of van Geemen and Rapagnetta \cite{vGR26},
  pull-back of the dual Beauville--Bogomolov class of a hyperk\"ahler
  manifold: up to multiples of $\pi_{0}$, the exceptional classes are exactly
  such twisted dual forms of the part of $H^{2}$ of the square generated by
  its symplectic class, and a rational map onto a
  symplectic subvariety of some $S_{\mu}^{[n]}$, $n\ge5$, would suffice;
  K3 surfaces and hyperk\"ahler eightfolds cannot serve
  (\Cref{prop:hksquare}, \Cref{rem:hkroute}).
\item \emph{Motivated classes.} Every Hodge class on an abelian variety is
  motivated \cite[Th\'eor\`eme~0.6.2]{And96b}, so $\omega_{t}$ is motivated for
  every $t$; a motivated class is algebraic once the Lefschetz involutions of
  the auxiliary varieties are (\Cref{prop:lefmot}). Along the family the
  auxiliary variety that suffices is $W\times_{C}W$
  (\Cref{thm:lefschetzpropagation}), and its Lefschetz standard conjecture is
  equivalent to the target (\Cref{cor:mumfordequiv}): the route moves the
  problem and does not shorten it.
\end{enumerate}
Five of the ten routes are closed off by theorems. Of the other five, two
work at special points only, the CM points and the closed points of the
reductions modulo $p$, and each suffices exactly when a degree bound holds;
two are reduced to a single object, one Kuga--Satake class or one perfect
complex with $\dim\Ext^{2}=119$ at one CM point; and one is equivalent to the
target. They are not five separate obligations: three of them are the target itself
in other forms and the other two are stronger inputs that imply it, so any
one of them would suffice (\Cref{cor:mumfordone}, \Cref{rem:mumfordneeded}).
\end{remark}

\begin{figure}[tbp]
\centering
  \includegraphics[width=0.96\textwidth]{fig_bypass}
\caption{\Cref{rem:bypassaudit}: the ten routes to the Mumford target that
avoid the Lefschetz standard conjecture. Clay, with a crossed dashed spoke: a
route closed off by a theorem of \Cref{ssec:mumfordrigid} or
\Cref{ssec:mumfordobject}. Grass: a route that works at special points only,
at the CM points or at the closed points of the reductions modulo $p$, and
suffices exactly when a degree bound holds (\Cref{cor:mumfordbounded},
\Cref{thm:modpdensity}). Indigo: a route reduced to one statement, of
the same kind as the target or equivalent to it.}
\label{fig:bypass}
\end{figure}
